\documentclass[twoside]{article}

\usepackage{aistats2027}
% Camera-ready: \usepackage[accepted]{aistats2027}
% Non-anonymous preprint: \usepackage[preprint]{aistats2027}

% natbib, author-year (required for camera-ready; allowed for submission)
\usepackage[round]{natbib}
\renewcommand{\bibname}{References}
\renewcommand{\bibsection}{\subsubsection*{\bibname}}

\usepackage[utf8]{inputenc}
\usepackage{amsmath}
\usepackage{amssymb}\usepackage{amsthm}
\usepackage{amsfonts}
\usepackage{mathtools}
\usepackage{mathrsfs}
\usepackage{bm}
\usepackage{graphicx}
\usepackage{subcaption}
\usepackage{url}
\usepackage{booktabs}
\usepackage{array}
\usepackage{nicefrac}
\usepackage{microtype}
\usepackage{physics}
\usepackage{wrapfig}
\usepackage{algorithm,algpseudocode}
\usepackage{xcolor}
\usepackage{enumitem}
\usepackage{hyperref}
\hypersetup{hidelinks, pdftitle={}, pdfauthor={}, pdfsubject={}, pdfkeywords={}}

%--- Blackboard-bold sets ---
\newcommand{\R}{\mathbb{R}}
\newcommand{\N}{\mathbb{N}}
\newcommand{\Z}{\mathbb{Z}}
\newcommand{\Q}{\mathbb{Q}}
\newcommand{\C}{\mathbb{C}}
\newcommand{\E}{\mathbb{E}}
\newcommand{\Tdist}{T}
\newcommand{\statespace}{\mathcal{S}}
\newcommand{\actionspace}{\mathcal{A}}
\newcommand{\obsspace}{\mathcal{O}}
\newcommand{\paramspace}{\Theta}

%--- Probability & expectation ---
\newcommand{\Prb}{\mathbb{P}}
\newcommand{\Var}{\mathrm{Var}}
\newcommand{\Cov}{\mathrm{Cov}}

%--- Common math operators ---
\newcommand{\diag}{\operatorname{diag}}
\newcommand{\supp}{\operatorname{supp}}
\newcommand{\argmax}{\operatorname*{arg\,max}}
\newcommand{\argmin}{\operatorname*{arg\,min}}
\newcommand{\ind}{\mathbf{1}}
\newcommand{\inner}[2]{\left\langle #1, #2 \right\rangle}

%--- Calligraphic sets ---
\newcommand{\cA}{\mathcal{A}}
\newcommand{\cS}{\mathcal{S}}
\newcommand{\cT}{\mathcal{T}}
\newcommand{\cF}{\mathcal{F}}
\newcommand{\cE}{\mathcal{E}}

%--- RL / MDP notation ---
\newcommand{\SSS}{\mathcal{S}}
\newcommand{\AAA}{\mathcal{A}}
\newcommand{\PP}{\mathcal{P}}
\newcommand{\RRR}{\mathcal{R}}
\newcommand{\gammaRL}{\gamma}

%--- Confidence bounds ---
\newcommand{\olQ}{\overline{Q}}
\newcommand{\ulQ}{\underline{Q}}
\newcommand{\olV}{\overline{V}}
\newcommand{\ulV}{\underline{V}}

%--- Miscellaneous ---
\newcommand{\eps}{\varepsilon}
\newcommand{\dt}{\mathrm{d}t}
\newcommand{\given}{\,\bigl\vert\,}
\newcommand{\half}{\tfrac12}
\newcommand{\thalf}{\tfrac32}

%--- Author notes (remove before submission) ---
\newcommand{\todo}[1]{\textcolor{red}{[TODO: #1]}}

\renewenvironment{proof}{\paragraph{Proof:}}{\hfill$\square$}
\newenvironment{solution}{\paragraph{Solution:}}{\hfill$\square$}

\newtheorem{theorem}{Theorem}
\newtheorem{proposition}[theorem]{Proposition}
\newtheorem{lemma}[theorem]{Lemma}
\newtheorem{corollary}[theorem]{Corollary}
\newtheorem{definition}[theorem]{Definition}
\newtheorem{assumption}[theorem]{Assumption}
\newtheorem{remark}[theorem]{Remark}
\newtheorem{example}[theorem]{Example}
\newtheorem{conjecture}[theorem]{Conjecture}
\newtheorem{hypothesis}[theorem]{Hypothesis}

\begin{document}

% Used only in camera-ready mode (the full title is too long for the running head).
\runningtitle{Predicting Level Difficulty Before Running It}
% \runningauthor{Chang}  % camera-ready only

\twocolumn[
\aistatstitle{Predicting Level Difficulty Before Running It:\\ Reliable Labels, Layout Graders, and Cheaper Curricula}

\aistatsauthor{ William Chang }

\aistatsaddress{ Department of Applied Mathematics, University of California, Los Angeles } ]

\begin{abstract}
Curriculum methods for reinforcement learning choose training levels by running the agent inside each candidate, so scoring levels is a major cost. We ask how much of a level's difficulty can be read from its layout before anything is run. First, we show that the difficulty labels such predictions are judged against are often unreliable, and we require every prediction claim to pass a split-half reliability check. Second, guided by a bound that ties changes in a level's dynamics to changes in its graph of feasible moves, we build a grader from layout, reachability, and Laplacian-spectrum features that scores a level in milliseconds. On spatial domains (Sokoban and two MiniGrid lava tasks) it is the strongest execution-free predictor we tested, can be worth many rollouts of a trained agent, and needs fewer labels than a CNN; on logic-heavy Karel tasks it fails. Third, used as a curriculum signal, it gives an easy-to-hard order that beats a random order in a pre-registered test, ranks levels within a grid size, and matches the generator's own difficulty setting without being told it; inside SFL it matches the standard method with about a quarter of the scoring rollouts. The canonical curriculum testbeds are spatial, the regime the method covers.
\end{abstract}

\section{Introduction}
\label{sec:intro}

Automated curricula train reinforcement learning agents on progressively harder tasks supplied by an environment generator---evolutionary level editors \citep{wang2019poet, parker2022evolving}, learned adversaries \citep{dennis2020emergent}, replay-based curators over procedural samplers \citep{jiang2021prioritized, jiang2021replay}, or language models that write environment programs \citep{ma2024eureka, liang2024eurekaverse, faldor2025omni}. However they generate, they share one loop: propose a candidate environment, score its usefulness to the current agent, train on it if it scores well.

Scoring dominates that loop's compute budget, because every established scoring function is defined on trajectories the agent produces \emph{inside} the candidate: temporal-difference error \citep{jiang2021prioritized}, approximate regret \citep{dennis2020emergent, parker2022evolving}, transition-prediction error \citep{cho2025traced}, binary-outcome learnability $p(1-p)$ \citep{tzannetos2023proximal, rutherford2024no}. This is affordable when the generator is a fast procedural sampler, and far less so when each candidate costs minutes to instantiate and the pool is large. The asymmetry motivates a question the curriculum literature has largely not asked: \textbf{how much of an environment's difficulty is legible before any policy is executed in it?}

\paragraph{Why existing tools do not answer it.} When environments share dynamics and differ only in reward, successor features with generalized policy improvement already answer it from the specification alone \citep{barreto2017successor, barreto2018transfer}. A generator that edits a layout, however, changes the dynamics, and the bounds that cover changed dynamics need objects unavailable before deployment: the simulation lemma and bisimulation metrics need the transition kernels \citep{kearns2002near, ferns2004metrics}, and successor-feature extensions need target-environment samples \citep{abdolshah2021new, malekzadeh2023uncertainty}.

\paragraph{The structural observation.} A large class of generated environments---grid worlds, tile-based levels, layout-parameterized simulators---differ mainly in \emph{which one-step transitions are feasible}. Feasibility is a combinatorial object, the feasibility graph $G_\theta$, that the generator has already written down. Theorem~\ref{prop:adjacency-control} shows that, for deterministic dynamics on a shared grid, the change in a policy's transition operator between two levels is bounded by the number of feasible moves that changed. This points to features of $G_\theta$---reachability, bottlenecks, and the Laplacian spectrum---as the place to look for difficulty. The bound is loose; it tells us which features to compute, and the experiments test whether they work.

\paragraph{The target must be measurable before it can be predicted.} Curriculum work routinely treats measured agent performance on a level as its difficulty without checking that the measurement repeats. It often does not. A single hill-climbing run is $1$--$35\%$ signal across twelve Karel tasks, and a reference PPO oracle used in earlier curriculum work gave $91\%$ of its levels the same label. We therefore check every label with a split-half reliability statistic (Eq.~\ref{eq:reliability}) before predicting it, and use $\sqrt{\rho_{\mathrm{rel}}}$ as the best score any predictor could reach.

\paragraph{Contributions.}
\begin{enumerate}[leftmargin=*, itemsep=0pt, topsep=2pt]
\item \textbf{Measurement.} A reliability gate for difficulty labels, with variance components showing how many search runs each task needs (from $4$ to over $140$), and evidence that the gate also catches degenerate oracles (Section~\ref{sec:reliability}).
\item \textbf{A layout grader with a reason to work.} A bound on the change in dynamics by the distance between feasibility graphs (Theorem~\ref{prop:adjacency-control}) motivates a grader from layout, reachability, and Laplacian-spectrum features, and a layer ablation confirms that the graph features carry the signal (Sections~\ref{sec:prelim} and~\ref{sec:results-prediction}).
\item \textbf{Prediction results on identical levels.} On spatial domains the grader beats every execution-free baseline we tested, matches $16$ trained-agent rollouts on the hardest pool, and on Sokoban needs fewer labels than a CNN; on logic-heavy Karel it fails (Section~\ref{sec:results-prediction}).
\item \textbf{A cheap curriculum signal.} In pre-registered tests the grader's order beats a random order and matches the generator's own difficulty setting without being told it, and it ranks levels within a grid size rather than recovering size; inside SFL \citep{rutherford2024no} it matches the standard version at about a quarter of the scoring rollouts (Section~\ref{sec:results-curriculum-v2}).
\end{enumerate}

\paragraph{Scope: how much of the field is spatial?}
\label{sec:results-scope}
\label{sec:results-cross}
The method targets \emph{spatial} difficulty: levels whose generator changes the layout while the rules stay fixed. To judge how often this applies, we checked what the level generator varies in the main curriculum and environment-design benchmarks, against each paper's text (Table~\ref{tab:survey}; quotes in Appendix~\ref{app:survey}). The canonical testbed of the line PAIRED $\to$ ACCEL $\to$ SFL $\to$ JaxUED is a MiniGrid maze whose generator moves only walls, goal, and agent, which is exactly the spatial regime. Logical generators exist and produce tasks at scale: XLand-MiniGrid samples from a bank of rule sets, and our method does not cover it. Procgen's reactive games, whose difficulty lies in spawn timing, and physics-based generators are also outside what a static layout can show.

\begin{table}[t]
\centering
\small
\caption{What the level generator varies in common curriculum benchmarks, verified against each paper.}
\label{tab:survey}
\setlength{\tabcolsep}{3pt}
\begin{tabular}{p{2.9cm}p{2.7cm}l}
\toprule
Benchmark (used by) & Generator varies & Category \\
\midrule
MiniGrid maze (PAIRED, ACCEL, SFL, JaxUED) & walls, goal, start & spatial \\
MultiRoom (PLR) & rooms, positions & spatial \\
JaxNav (SFL) & obstacle map, start, goal & spatial \\
MiniHack lava (ACCEL) & lava tiles & spatial \\
Procgen (PLR) & layout, entities, timing & mostly spatial \\
ObstructedMaze (PLR) & layout, key chain & mixed \\
BabyAI & instruction, layout & mixed \\
Craftax & terrain; fixed tech tree & mixed \\
XLand-MiniGrid (SFL) & rule sets, goals & logical \\
BipedalWalker, Kinetix & terrain, physics & out of scope \\
\bottomrule
\end{tabular}
\end{table}

\section{Related Work}
\label{sec:related-main}

\paragraph{Curriculum scoring and its cost.} Unsupervised environment design scores levels by regret or value error \citep{dennis2020emergent, jiang2021prioritized, parker2022evolving}, and SFL by the learnability $p(1-p)$ of rollouts \citep{rutherford2024no}. Methods that cut this cost still run the environment, through a provisional policy update \citep{yuan2026pace}, sampled episodes \citep{furuta2021policy}, or an agent population whose performance embeds tasks \citep{mahajan2024learning}. Our grader runs nothing at use time.

\paragraph{Predicting difficulty.} Static-feature difficulty estimation is long established for Sokoban \citep{ashlock2010evolution, vankreveld2015automated, kartal2016data}, and regressing solver performance on instance features is an empirical hardness model \citep{leytonbrown2009empirical, hutter2014algorithm}. Closest to us, \citet{krsteski2026predicting} predict agentic-task difficulty before execution (Spearman $0.40$ in distribution, $0.23$ on new benchmarks) but do not check label reliability. We add the reliability gate, a value-gap argument for the features, head-to-head comparisons on identical levels, and a curriculum test (extended discussion in Appendix~\ref{app:related}).

\section{Setting and the Feasibility Graph}
\label{sec:prelim}
\label{sec:dynamics-bound}

\paragraph{Setting.} A generator outputs parameters $\theta$, each defining a Markov decision process $\mathcal{M}_\theta = (\statespace, \actionspace, \Tdist_\theta, R_\theta, I, \gamma)$ whose states, actions, start distribution, and discount are \emph{shared} across the family; only the transition kernel $\Tdist_\theta$ and reward $R_\theta$ change. For a policy $\pi$, write $T^\pi_\theta(u,v) = \sum_a \pi(a\mid u)\Tdist_\theta(v \mid u,a)$ for the transition operator it induces. Given levels with measured difficulty, we want to predict the difficulty of a new level $\theta_{\text{new}}$ without running anything in it. When levels differ only in reward, successor features already do this from the specification \citep{barreto2017successor, barreto2018transfer}; a generator that edits a layout changes $\Tdist_\theta$, and existing bounds for changed dynamics need the kernels or target samples (Appendix~\ref{app:valuegap} gives the value-gap decomposition and how our bound enters it).

\begin{definition}[Feasibility graph]
\label{def:kinograph}
Let $\mathcal{V} \subset \statespace$ be a discretization of the shared state space (e.g.\ the cells of a grid). The feasibility graph of $\theta$ is $G_\theta = (\mathcal{V}, \mathcal{E}_\theta)$ with $(u,v) \in \mathcal{E}_\theta$ if and only if a one-step move from $u$ to $v$ is possible under $\theta$. Write $A_\theta$ for its adjacency matrix and $D_\theta$ for its degree matrix.
\end{definition}

All levels of a family share the vertex set and differ only in edges, so $\|A_\theta - A_{\theta'}\|_F^2$ counts the moves whose feasibility changed. The graph is unweighted: it records whether a move is possible, not how likely (Remark~\ref{rem:weights}).

\begin{theorem}[Adjacency control of the transition operator]
\label{prop:adjacency-control}
Let $G_\theta, G_{\theta'}$ be feasibility graphs on a finite shared vertex set $\mathcal V$ with finite action set $\mathcal A$, and fix a stochastic policy $\pi$. Assume \textbf{support consistency}: $\Tdist_\theta(v\mid u,a)>0 \Rightarrow A_\theta(u,v)=1$, and likewise for $\theta'$, so each kernel places mass only on feasible one-step moves. Let $d_{\max} := \max_{u} \deg_{G_\theta\cup G_{\theta'}}(u)$ and suppose at most $d_{\max}$ actions send $u$ to $v$ in either environment. Define the \emph{shared-edge residual}, measuring how far the two environments disagree on edges feasible in \emph{both},
\begin{equation}\label{eq:residual}
\begin{split}
    \Delta_\pi(\theta,\theta')
    :=
    \Bigg(
    &\sum_{\substack{u,v:\\ A_\theta(u,v)=A_{\theta'}(u,v)=1}}
    \Big|
        \sum_{a}
        \pi(a\mid u)
        \\
        &\times\big(
            \Tdist_\theta(v\mid u,a) - \Tdist_{\theta'}(v\mid u,a)
        \big)
    \Big|^{2}
    \Bigg)^{1/2}.
\end{split}
\end{equation}
Then, with $\|\pi\|_\infty := \max_{u,a}\pi(a\mid u)$ and $C_\pi := d_{\max}\|\pi\|_\infty \le d_{\max}$,
\begin{equation}\label{eq:adjacency-control}
    \big\|T^\pi_\theta - T^\pi_{\theta'}\big\|_{\mathrm{op}}
    \;\le\;
    C_\pi \big\|A_\theta-A_{\theta'}\big\|_F
    \;+\;
    \Delta_\pi(\theta,\theta').
\end{equation}
\end{theorem}

\noindent \emph{Proof idea} (Appendix~\ref{app:proofs}). Sort the entries of $T^\pi_\theta - T^\pi_{\theta'}$ by how feasibility compares across the two levels: pairs infeasible in both are zero, pairs feasible in exactly one are charged to $\|A_\theta - A_{\theta'}\|_F$, and pairs feasible in both form $\Delta_\pi$.

\begin{corollary}[Common-edge invariance]
\label{cor:common-edge}
If both levels move the agent identically along every edge feasible in both, then $\Delta_\pi=0$ and $\|T^\pi_\theta - T^\pi_{\theta'}\|_{\mathrm{op}} \le C_\pi \|A_\theta-A_{\theta'}\|_F$. This holds for deterministic dynamics on a shared grid, where $\theta$ only blocks moves (``go north from $u$'' lands on the same cell whenever it is allowed), and so for every domain in this paper.
\end{corollary}

\paragraph{What the theorem buys, and what it does not.} Combined with the standard resolvent argument, the corollary bounds the dynamics part of the value gap between two levels by a constant times $\|A_\theta-A_{\theta'}\|_F$ (Eq.~\ref{eq:final-bound}, Appendix~\ref{app:valuegap}): for a fixed policy, levels whose feasibility graphs are close have close values. Difficulty should therefore be visible in how $G_\theta$ is connected, which motivates the grader's features in Section~\ref{sec:geometric-predictor}. Two caveats keep this honest. The constant exceeds $10^6$ at our grid sizes (Remark~\ref{rem:ellinfty}), so the bound guides the choice of features and predicts no number. And the grader scores one level at a time with a learned model rather than comparing pairs of graphs, so whether those features predict difficulty is an empirical question; the layer ablation of Section~\ref{sec:results-prediction} is our evidence that they do.

\section{Measuring Difficulty Before Predicting It}
\label{sec:framework}
\label{sec:reliability}

\paragraph{Difficulty labels.}
\label{sec:difficulty-label}
A label turns ``how hard is level $\theta$'' into a number by running an expensive reference solver on it. For search-labelled domains we use hill climbing over a DSL program space \citep{trivedi2021learning, carvalho2024reclaiming}; with $r^\star_{\mathrm{HC}}(\theta;\xi)\in[0,1]$ the best task-normalized reward under search seed $\xi$, the label averages $R$ seeds,
\begin{equation}
    D_{\mathrm{HC}}(\theta) \;=\; 1 - \frac{1}{R}\sum_{r=1}^{R} r^\star_{\mathrm{HC}}(\theta; \xi_r).
    \label{eq:difficulty-label}
\end{equation}
For agent-labelled domains the label is $1-$success of independently trained reference PPO agents, and for the public Sokoban set we use labels published by others (Section~\ref{sec:setup}). Under a sparse reward almost every program scores zero and (\ref{eq:difficulty-label}) has no variance, so for navigation we shape the reward by shortest-path progress, $r_{\mathrm{shaped}} = (d_{\mathrm{start}} - d_{\min})/d_{\mathrm{start}}$ (Appendix~\ref{app:method}).
\label{sec:shaping}

\paragraph{The reliability gate.} A low predictor score means nothing if re-measuring the label reorders the levels. We split the $R$ runs into halves, recompute the label in each, and correct their rank agreement to full length (Spearman--Brown), averaging over $20$ random splits:
\begin{equation}
    \rho_{\mathrm{rel}} \;=\; \E_{\text{splits}}\!\left[\frac{2\,\rho_s(D^{(A)}, D^{(B)})}{1 + \rho_s(D^{(A)}, D^{(B)})}\right].
    \label{eq:reliability}
\end{equation}
Because the label is a noisy copy of true difficulty, no predictor can exceed a Spearman of about $\sqrt{\rho_{\mathrm{rel}}}$ against it; we report this \emph{ceiling} with every result.

\paragraph{Finding 1: single runs are mostly noise.} A one-way variance decomposition (grid versus search seed) on the twelve Karel tasks shows that one search run is $1$--$35\%$ signal (Table~\ref{tab:karel}). The number of runs needed for $\rho_{\mathrm{rel}}=0.7$ ranges from $4$ (\textsc{WallAvoider}) to over $140$ (\textsc{Seeder}, \textsc{Harvester}). With $R=12$, five tasks pass the gate; the rest are reported but not used for prediction claims.

\begin{table}[t]
\centering
\small
\caption{How measurable is Karel difficulty? Single-run share of signal (intraclass correlation), split-half reliability at $R=12$, and runs needed for reliability $0.7$. Tasks above the line pass the gate.}
\label{tab:karel}
\setlength{\tabcolsep}{4pt}
\begin{tabular}{lccc}
\toprule
Task & Single run & $\rho_{\mathrm{rel}}$ ($R{=}12$) & Runs for $0.7$ \\
\midrule
\textsc{WallAvoider} & $0.35$ & $0.82$ & $4$ \\
\textsc{DoorKey}     & $0.18$ & $0.71$ & $10$ \\
\textsc{Snake}       & $0.17$ & $0.71$ & $11$ \\
\textsc{Maze}        & $0.10$ & $0.70$ & $21$ \\
\textsc{TopOff}      & $0.19$ & $0.57$ & $10$ \\
\midrule
\textsc{PathFollow}  & $0.18$ & $0.43$ & $10$ \\
\textsc{OneStroke}   & $0.05$ & $0.44$ & $42$ \\
\textsc{StairClimber}& $0.04$ & $0.25$ & $54$ \\
\textsc{CleanHouse}  & $0.03$ & $0.24$ & $67$ \\
\textsc{FourCorners} & $0.03$ & $0.23$ & $79$ \\
\textsc{Seeder}      & $0.02$ & ${\approx}0$ & $144$ \\
\textsc{Harvester}   & $0.01$ & ${\approx}0.05$ & $254$ \\
\bottomrule
\end{tabular}
\end{table}

\paragraph{Finding 2: the gate catches bad oracles.} A reference-PPO oracle used in earlier curriculum work on \textsc{LavaCrossing} gave $82$ of $90$ levels the label $1.0$. Rebuilding it with three fresh reference agents still left $76\%$ of levels tied on the large grids, with within-cell reliability $0.21$: an agent trained without a curriculum cannot label the levels where curricula are needed. A progress-based label was reliable there ($0.83$) but almost fully explained by where the first lava gap sits on the path ($\rho_s=-0.82$), so reliability is necessary but not sufficient.

\section{The Grader and the Baselines}
\label{sec:geometric-predictor}

\paragraph{The grader.}
\label{sec:geometric-descriptors}
From the layout alone we build $G_\theta$ over free cells and compute three layers of features, $g(\theta) = [g_{\mathrm{loc}};\, g_{\mathrm{reach}};\, g_{\mathrm{spec}}]$. \emph{Local} counts (free cells, hazard fraction, dead-ends, corridor and branch cells) cannot tell whether a path exists. \emph{Reachability} features of $G_\theta$ (start eccentricity, start--goal distance and its excess over Manhattan distance, bridges, articulation points) can. \emph{Spectral} features, the eigenvalues $\lambda_2,\ldots,\lambda_6$ of the normalized Laplacian $L_\theta = I - D_\theta^{-1/2}A_\theta D_\theta^{-1/2}$, summarize bottlenecks globally: by the Cheeger inequality \citep{chung1997spectral}, a small $\lambda_2$ means the free region is nearly cut in two.
\begin{equation}
    g(\theta) \;=\; \big[\, g_{\mathrm{loc}}(\theta) \;;\; g_{\mathrm{reach}}(\theta) \;;\; g_{\mathrm{spec}}(\theta) \,\big] \;\in\; \R^{k}.
\label{eq:full-descriptor}
\end{equation}
A random forest ($300$ trees) maps $g(\theta)$ to difficulty. Scoring a level takes milliseconds and runs nothing in it. The forest is fit once on labelled levels and scored out-of-fold ($5$-fold) or on a separate pool, so its labels are a one-time cost per generator. It is never given grid size or river count, although the free-cell count carries size.

\paragraph{Baselines, all on the same levels and labels.} (i) \textbf{Kartal 2016}, a published unfitted Sokoban difficulty formula \citep{kartal2016data}. (ii) \textbf{Designer knobs}: generator settings that index difficulty, where they exist (grid size plus river count in \textsc{LavaCrossing}). (iii) \textbf{LLM judge}: Qwen2.5-7B-Instruct \citep{hui2024qwen2} shown the rules and the ASCII layout and asked for a $0$--$9$ difficulty, scored by the expected digit; one fixed prompt, no tuning. (iv) \textbf{CNN on the raw grid}: three convolutional layers on a one-hot grid, trained on the same folds with the eight rotations and flips of each level. (v) \textbf{Running the expensive thing}: $k$ episodes of a trained reference agent, or $k$ runs of the reference search, which is what rollout-based curriculum methods pay for.

\section{Experimental Setup}
\label{sec:setup}

We use four spatial domains and one logic-heavy domain (Table~\ref{tab:domains}).
\textbf{Sokoban, public} \citep{guez2019investigation}: $3{,}000$ \textsc{Boxoban} validation levels ($1{,}500$ unfiltered, $1{,}500$ medium), labelled by published A$^\star$ search effort and solution length \citep{taufeeque2024planning}, and by tier (medium versus unfiltered, AUC). These labels come from a deterministic solver and are not ours.
\textbf{Sokoban, $200$ levels}: our budgeted best-first search label, stratified over tiers.
\textbf{MiniGrid \textsc{LavaGap}}: one lava strip with a single gap; $60$ levels, shaped hill-climbing label.
\textbf{MiniGrid \textsc{LavaCrossing}}: $1$--$3$ lava rivers, each with one crossing; two pools of $90$ levels, grid sizes $\{7,9,11\}$ and $\{9,13,17\}$, labelled by $1-$success of three reference PPO agents ($600$k steps, $20$ episodes per level).
\textbf{Karel} \citep{pattis1981karel, trivedi2021learning}: twelve tasks, $80$--$200$ grids each, hill-climbing label (Table~\ref{tab:karel}).
We report out-of-fold Spearman $\rho$ against the label (AUC for tier), with the ceiling $\sqrt{\rho_{\mathrm{rel}}}$.

\begin{table}[t]
\centering
\small
\caption{Domains and labels. Ceiling $=\sqrt{\rho_{\mathrm{rel}}}$, the best Spearman any predictor can reach against the label.}
\label{tab:domains}
\setlength{\tabcolsep}{3pt}
\begin{tabular}{llcc}
\toprule
Domain & Label & Levels & Ceiling \\
\midrule
Sokoban public & A$^\star$ effort, length, tier & $3{,}000$ & $1.00$ \\
Sokoban & budgeted search & $200$ & $0.99$ \\
\textsc{LavaGap} & shaped hill climbing & $60$ & $0.89$ \\
\textsc{LavaCrossing} 7/9/11 & $1-$PPO success & $90$ & $0.75$ \\
\textsc{LavaCrossing} 9/13/17 & $1-$PPO success & $90$ & $0.76$ \\
Karel (5 gated tasks) & hill climbing & $80$--$200$ & $0.75$--$0.91$ \\
\bottomrule
\end{tabular}
\end{table}

\section{Results: Predicting Difficulty}
\label{sec:results}
\label{sec:results-prediction}
\label{sec:results-sokoban}

Table~\ref{tab:prediction} compares every method on identical levels and labels.

\begin{table*}[t]
\centering
\small
\caption{Spearman $\rho$ between each method's score and the label (AUC for tier), on identical levels. Ours, Kartal, knobs and the LLM execute nothing; the CNN is a learned model on the raw grid (run where enough labelled levels exist); rollouts are $k$ agent runs, i.e.\ $k$ episodes of a trained reference PPO agent on every level. Ceiling $=\sqrt{\rho_{\mathrm{rel}}}$. n/a: the method does not exist for that domain (Kartal is Sokoban-only; \textsc{LavaGap} has no difficulty knob; on public Sokoban the tier is a label, and the labels come from a solver, not an agent). $^\ast$The Boxoban tier is computed by running agents, so it is not a free knob. Bootstrap $95\%$ intervals for ours (resampling levels): public Sokoban $\pm0.01$--$0.02$; Sokoban $200$ $[0.20, 0.44]$; \textsc{LavaGap} $[0.38, 0.70]$; \textsc{LavaCrossing} $[0.25, 0.52]$ and $[0.29, 0.60]$.}
\label{tab:prediction}
\setlength{\tabcolsep}{5pt}
\begin{tabular}{llccccccc}
\toprule
Domain & Label & Ceiling & \textbf{Ours} & Kartal & LLM & Knobs & CNN & Rollouts $k{=}1 \to 16$ \\
\midrule
Sokoban public ($3{,}000$) & A$^\star$ search effort & $1.00$ & $\mathbf{0.75}$ & $0.41$ & $0.18$ & n/a & $0.77$ & n/a \\
 & A$^\star$ solution length & $1.00$ & $0.71$ & $0.48$ & -- & n/a & $0.77$ & n/a \\
 & tier (AUC) & -- & $0.80$ & $0.64$ & $0.57$ & n/a & $0.82$ & n/a \\
Sokoban ($200$) & budgeted search & $0.99$ & $0.32$ & n/a & $0.01$ & $0.57^\ast$ & -- & n/a \\
\textsc{LavaGap} ($60$) & shaped search & $0.89$ & $\mathbf{0.57}$ & n/a & $0.29$ & n/a & -- & see Table~\ref{tab:worth} \\
\textsc{LavaCrossing} 7/9/11 ($90$) & $1-$PPO success & $0.75$ & $0.40$ & n/a & $0.20$ & $0.24$ & -- & $0.42 \to 0.49$ \\
\textsc{LavaCrossing} 9/13/17 ($90$) & $1-$PPO success & $0.76$ & $\mathbf{0.46}$ & n/a & $0.27$ & $0.35$ & -- & $0.36 \to 0.44$ \\
\bottomrule
\end{tabular}
\end{table*}

\paragraph{The grader is the strongest execution-free method on spatial domains.} On public Sokoban it beats the published Kartal formula by $0.23$--$0.34$ and the LLM judge by more. On \textsc{LavaCrossing} 9/13/17 it was fit only on the separate 7/9/11 pool, yet it beats the designer knobs ($0.46$ vs $0.35$) and matches $16$ rollouts of a trained agent ($0.44$). On \textsc{LavaGap} it reaches $0.57$ of a $0.89$ ceiling ($0.56 \pm 0.04$ over $10$ cross-validation splits).

\paragraph{Where it loses, and why.} Three cells go against us. (i) With all $3{,}000$ labels, a CNN on the raw grid is slightly better ($0.77$ vs $0.75$); Section~\ref{sec:results-worth} shows this reverses when labels are scarce. (ii) On the $200$-level Sokoban pool the Boxoban tier beats us ($0.57$ vs $0.32$), but that tier was made by running agents. (iii) On \textsc{LavaGap}, hand-designed gap-position features reach $0.77$; we report them as a hand-engineered reference, not as our method.

\paragraph{Secondary: grader plus LLM.} Adding the LLM score as one more input raises \textsc{LavaGap} to $0.82$ ($0.82\pm0.02$ over $10$ splits) and helps on some Karel tasks, but gives nothing on Sokoban or \textsc{LavaCrossing}. Because it needs an LLM call per level, we treat it as a secondary result (Appendix~\ref{app:extra}).

\paragraph{Which feature layers matter?} Table~\ref{tab:ablation} gives the grader only some of its three layers (Section~\ref{sec:geometric-descriptors}). On every spatial domain, going beyond local counts helps, by $+0.06$ to $+0.31$. The spectral layer adds on \textsc{LavaCrossing} ($+0.06$) and Sokoban ($+0.01$--$0.02$, well outside the $\pm0.003$ split-to-split spread) but not on \textsc{LavaGap}, whose single bottleneck the reachability features already capture. The path and connectivity quantities that the bound charges therefore carry signal that simple counts do not.

\begin{table}[t]
\centering
\small
\caption{Feature-layer ablation: Spearman (AUC for tier) of the same forest given only some layers; mean over $10$ cross-validation splits ($3$ for public Sokoban). Both \textsc{LavaCrossing} rows use within-pool cross-validation, so they differ from Table~\ref{tab:prediction}, where the 9/13/17 grader is fit on the separate 7/9/11 pool. $^\dagger$Adding the Sokoban-specific box and target features gives $0.75$ and $0.80$.}
\label{tab:ablation}
\setlength{\tabcolsep}{3pt}
\begin{tabular}{lccc}
\toprule
Domain & Local & $+$Reach & $+$Spectral \\
\midrule
\textsc{LavaGap} & $0.26$ & $\mathbf{0.57}$ & $0.56$ \\
\textsc{LavaCrossing} 7/9/11 & $0.27$ & $0.34$ & $\mathbf{0.39}$ \\
\textsc{LavaCrossing} 9/13/17 & $0.26$ & $0.28$ & $\mathbf{0.34}$ \\
Sokoban, A$^\star$ effort & $0.59$ & $0.66$ & $\mathbf{0.67}^\dagger$ \\
Sokoban, tier (AUC) & $0.74$ & $0.79$ & $\mathbf{0.80}^\dagger$ \\
\bottomrule
\end{tabular}
\end{table}

\subsection{What Is a Free Prediction Worth?}
\label{sec:results-worth}

\paragraph{In rollouts.} Table~\ref{tab:worth} (top) and Figure~\ref{fig:worth} score the grader against the alternative it replaces: running the expensive reference $k$ times on every level and using the observed failure rate. The reference is whatever produced that domain's label. On \textsc{LavaCrossing} it is an \emph{agent run}: one episode of a trained PPO agent, i.e.\ one play-through of the level. On \textsc{LavaGap} and Karel it is a \emph{search run}: one complete hill-climbing search of up to $2{,}000$ (\textsc{LavaGap}) or $5{,}000$ (Karel) candidate programs, so a single search run costs thousands of environment evaluations. To keep the comparison fair, the label is built from reference runs the estimator never sees. On the large \textsc{LavaCrossing} grids the free grader matches $16$ agent runs; on the small grids it is worth about one agent run, and on \textsc{LavaGap} about two search runs, i.e.\ roughly $4{,}000$ program evaluations per level.

\paragraph{In labels.} The grader is learned, so it needs labelled levels. Table~\ref{tab:worth} (bottom) trains the grader and the CNN on $n$ public Sokoban levels and tests both on the same $1{,}000$ held-out levels. The grader wins up to $n=500$, ties at $1{,}000$, and loses by $0.02$ at $2{,}000$. At $n=60$ it is also far more stable (tier AUC $0.75\pm0.01$ vs $0.66\pm0.11$ over five draws). Since each label costs a full solve, needing fewer of them is the practical advantage. We ran this comparison only on Sokoban, the one domain with enough labelled levels to train a CNN.

\begin{figure}[t]
\centering
\includegraphics[width=\columnwidth]{worth_curves.pdf}
\caption{What a free prediction is worth. \textbf{(a)} Each colour is one domain; compare a dashed line only with the solid line of the same colour. Dashed: our grader, which runs nothing. Solid: the failure rate after $k$ runs of that domain's expensive reference (agent runs for \textsc{LavaCrossing}, search runs for \textsc{LavaGap}). Where the dashed line lies above a solid point, the grader beats that many runs. \textbf{(b)} Public Sokoban: our grader and a CNN on the raw grid trained on $n$ labelled levels, scored on the same $1{,}000$ held-out levels (mean $\pm$ s.d.\ over $5$ draws).}
\label{fig:worth}
\end{figure}

\begin{table}[t]
\centering
\small
\caption{\textbf{Top:} the free grader versus $k$ runs of the expensive reference. \textsc{LavaCrossing}: agent runs (one episode of a trained PPO agent each). \textsc{LavaGap} and Karel: search runs (one full hill-climbing search each). To keep the comparison fair, the label here uses only $6$ held-out runs that no estimator sees, so ``Ours'' differs slightly from Table~\ref{tab:prediction} and Section~\ref{sec:results-karel}. \textbf{Bottom:} Spearman on $1{,}000$ held-out public Sokoban levels (A$^\star$ effort) after training on $n$ labels; mean of $5$ draws.}
\label{tab:worth}
\label{tab:fewlabels}
\setlength{\tabcolsep}{3.5pt}
\begin{tabular}{lccccc}
\toprule
 & Ours & $k{=}1$ & $k{=}2$ & $k{=}4$ & largest $k$ \\
\midrule
\textsc{LavaCross.} 9/13/17 & $\mathbf{0.46}$ & $0.36$ & $0.37$ & $0.40$ & $0.44$ ($16$) \\
\textsc{LavaCross.} 7/9/11  & $0.40$ & $0.42$ & $0.44$ & $0.46$ & $0.49$ ($16$) \\
\textsc{LavaGap}            & $0.49$ & $0.38$ & $0.49$ & $0.60$ & $0.65$ ($6$) \\
Karel \textsc{WallAvoider}  & $0.30$ & $0.53$ & $0.59$ & $0.65$ & $0.69$ ($6$) \\
Karel \textsc{DoorKey}      & $-0.09$ & $0.29$ & $0.41$ & $0.50$ & $0.55$ ($6$) \\
\midrule
Labels $n$ & $60$ & $200$ & $500$ & $1{,}000$ & $2{,}000$ \\
\midrule
Ours (forest) & $\mathbf{0.68}$ & $\mathbf{0.71}$ & $\mathbf{0.73}$ & $0.75$ & $0.75$ \\
CNN, raw grid & $0.61$ & $0.66$ & $0.69$ & $0.74$ & $0.77$ \\
\bottomrule
\end{tabular}
\end{table}

\subsection{Where the Grader Fails: Logical Difficulty}
\label{sec:results-karel}

Karel is the logic-heavy case: the layout changes, but what makes a grid hard is the program the agent must find. On the five tasks that pass the gate, the grader predicts \textsc{WallAvoider} ($0.36$) and partly \textsc{TopOff} ($0.28$) and nothing else (\textsc{DoorKey} $-0.14$, \textsc{Maze} $0.03$, \textsc{Snake} $0.07$). The LLM judge does better on those three ($0.41$, $0.17$, $0.19$) but no method exceeds $0.41$. Table~\ref{tab:worth} makes the failure concrete: on Karel the grader is worse than a \emph{single} search run, even on \textsc{WallAvoider}. The exception fits the story: \textsc{WallAvoider} asks the agent to fill the interior while avoiding the ring of wall-adjacent cells, so its difficulty is spatial. We conclude that the grader should not be used where difficulty comes from program logic.

\section{Results: A Cheap Curriculum Signal}
\label{sec:results-curriculum-v2}
\label{sec:results-curriculum-neural}

We do not ask whether curricula help; that is established. We ask how good an easy-to-hard order the grader gives, compared with other ways of getting one, and at what cost.

\paragraph{Setup (pre-registered).} We train PPO on \textsc{LavaCrossing} 9/13/17 ($90$ training and $36$ held-out levels, $1.5$M steps), where training without a curriculum fails. A shared controller splits the training levels into three bins by a difficulty score and unlocks the next bin when held success reaches $0.55$ or after $25\%$ of the budget; in practice unlocks are mostly deadline-driven. Conditions differ \emph{only} in the score: \textbf{flat} (no curriculum), \textbf{random order}, \textbf{rivers only} and \textbf{size order} (the two designer knobs), and the \textbf{grader}, fit only on the separate 7/9/11 pool. The metric, seeds $0$--$19$, and three one-sided Mann--Whitney tests (size order $>$ random, grader $>$ random, random $>$ flat) were fixed in writing before any run; seeds $20$--$39$ were a pre-specified secondary analysis.

\begin{table}[t]
\centering
\small
\caption{Curriculum on \textsc{LavaCrossing} 9/13/17: final held-out success (mean, $95\%$ CI) and seeds reaching $0.30$. Primary analysis $20$ seeds; secondary $40$.}
\label{tab:curriculum-v2}
\setlength{\tabcolsep}{3pt}
\begin{tabular}{lccc}
\toprule
Order & $20$ seeds & $40$ seeds & Reached \\
\midrule
Flat (none)   & $0.05$ & $0.06$ $[0.03,0.09]$ & $3/40$ \\
Random order  & $0.11$ & $0.12$ $[0.06,0.18]$ & $6/40$ \\
Rivers only   & $0.15$ & $0.11$ $[0.05,0.17]$ & $7/40$ \\
\textbf{Grader} & $\mathbf{0.31}$ & $\mathbf{0.35}$ $[0.25,0.45]$ & $23/40$ \\
Size order    & $0.35$ & $0.43$ $[0.35,0.51]$ & $30/40$ \\
\bottomrule
\end{tabular}
\end{table}

\paragraph{The grader gives a good order without being told the knob.} The grader beats a random order ($+0.19$, $p=0.016$ at $20$ seeds, still below $0.05$ after a Bonferroni correction over the three primary tests; $+0.23$, $p<0.001$ at $40$), and a random-order curriculum is no better than flat training ($p=0.80$), so the gain comes from the order itself (Table~\ref{tab:curriculum-v2}). The grader matches ordering by grid size, the generator's own difficulty setting ($-0.04$, $p=0.71$; $-0.08$, $p=0.35$ at $40$ seeds; size is numerically ahead). Grid size exists only because this generator's designer built difficulty along it, and choosing it requires knowing that in advance; river count, an equally simple setting, was no better than random. The grader needs no difficulty setting, so it also applies to generators such as Sokoban and \textsc{LavaGap} that have none, and it is computed once, before training.

\paragraph{Is the grader grid size in disguise?} No. Its free-cell feature carries size, and its scores correlate with size ($\rho_s=0.53$), but it also ranks levels \emph{within} a size. Against the $3$-agent label, its Spearman inside each size is $0.71$, $0.26$, and $0.28$ for sizes $9$, $13$, and $17$ (on the two larger sizes $83$--$87\%$ of labels are tied at total failure, which caps any ranking), and $0.24$ pooled within size ($95\%$ CI $[0.12, 0.50]$, permutation $p=0.009$); river count manages $0.10$--$0.24$. Its curriculum bins mix sizes: the easiest bin holds $20$ size-$9$ and $10$ size-$13$ levels, where size order's holds only size-$9$ levels. A within-cell check of this kind was run before the curriculum experiments ($+0.38$, recorded in the pre-registration); the per-size numbers above were computed afterwards on the same label.

\paragraph{The grader makes SFL cheaper.} SFL \citep{rutherford2024no} trains on levels the agent solves sometimes, found by rolling out every candidate. We compare standard SFL ($100$ candidates scored by rollouts), a cheap SFL scoring only $20$, and SFL with the grader: $20$ rollouts calibrate a map from grader score to success, and the grader then picks from all $100$. With the pre-registered $10$ seeds (Table~\ref{tab:sfl}) the grader version scores $0.25$ against $0.07$ for standard SFL while using $23\%$ of its scoring steps. The bootstrap interval for the difference, $[-0.05, 0.42]$, excludes losses larger than $0.10$, which was our pre-set criterion for ``matches standard SFL''. Its lead over cheap SFL ($+0.20$, $p=0.12$) was not significant. Ten further seeds, added after seeing these results and therefore exploratory, give $0.31$ vs $0.04$ for cheap SFL ($p=0.002$) and vs $0.13$ for standard SFL ($p=0.057$). One caveat: standard SFL learned little in this setting, so matching it is a modest bar. The refresh logs show why. Early in training almost no candidate is ever solved (median candidate success $0.005$; mean below $0.05$ in $79\%$ of refreshes), so $p(1-p)$ is near zero everywhere and SFL's buffer is close to a random draw, with $33\%$ size-$9$ levels, the same share as the pool. With the grader, $58\%$ of the buffer is size-$9$ levels: when rollouts carry no signal, a prior on which levels are easy does.

\paragraph{What the saving does not count.} The $23\%$ counts scoring during training. The grader's labels came from six reference PPO agents of $600$k steps each on the separate 7/9/11 pool, about $3.6$M steps paid once per generator. That is more than one SFL run's scoring budget, so the saving pays off only when the grader is reused across runs, seeds, or generator settings, which is how it was used here.

\begin{table}[t]
\centering
\small
\caption{SFL with and without the grader on \textsc{LavaCrossing} 9/13/17: final held-out success and scoring environment steps. Seeds $10$--$19$ are post hoc.}
\label{tab:sfl}
\label{tab:curriculum-cost}
\setlength{\tabcolsep}{3pt}
\begin{tabular}{lccc}
\toprule
 & Seeds $0$--$9$ & Seeds $0$--$19$ & Scoring steps \\
\midrule
Standard SFL ($100$ scored) & $0.07$ & $0.13$ & $793$k \\
Cheap SFL ($20$ scored) & $0.06$ & $0.04$ & $169$k \\
\textbf{SFL + grader} & $\mathbf{0.25}$ & $\mathbf{0.31}$ & $180$k \\
\bottomrule
\end{tabular}
\end{table}

\paragraph{Choosing \emph{which} levels does not beat random.} Using the grader to \emph{select} a training subset never beat a random subset of the same size: not on \textsc{LavaGap} with hill climbing, not on Sokoban, and not on \textsc{LavaCrossing} with PPO (Appendix~\ref{sec:results-curriculum}). The grader's value for curricula lies in ordering and in replacing scoring rollouts, not in selection.

\paragraph{Where no order can help.} Ordering signals can only be compared where order matters. In two settings flat training already succeeds, and there no order, the designer's or ours, separated from any other. On the small 7/9/11 grids flat training and size order did not differ ($0.43$ vs $0.53$, $10$ seeds). We also pre-registered a second, knob-free test on structured $19\times19$ mazes (Appendix~\ref{app:maze}); flat training reached at least $0.25$ on $14$ of $20$ seeds, and flat ($0.40$), random order ($0.35$), wall count ($0.47$), path length ($0.43$), and the grader ($0.42$) were statistically indistinguishable (all $p\ge0.05$). Making the mazes harder by keeping fewer shortcuts overshot: in pre-registered pilots neither flat training nor the wall-count order learned (mean success $\le0.03$), so this generator offered no setting where a curriculum was both needed and sufficient, and the main test was not run.

\section{Conclusion and Limitations}
\label{sec:conclusion}

How much of a level's difficulty can be read before anything is run? For spatial levels, a useful amount: a millisecond grader is the strongest execution-free predictor we tested, matches $16$ agent rollouts on the hardest pool, needs fewer labels than a CNN on Sokoban, gives a curriculum order as good as the generator's own difficulty setting without being told it, and lets SFL match its standard version at about a quarter of the scoring cost. For logic-heavy levels it fails. None of this can be judged without first checking that the difficulty label is measurable, which many labels are not.

\paragraph{Limitations.}
\label{sec:limitations-main}
(i) The grader matches but does not beat ordering by grid size, and its curriculum value is shown on one generator only: on a knob-free maze generator we found no setting where a curriculum was both needed and sufficient (mazes were either solved without one or not learned even with one). (ii) The grader needs labels, a one-time cost per generator that exceeds one SFL run's scoring budget, and with $1{,}000$ or more labels a CNN catches up. (iii) It fails on logical difficulty, and some popular benchmarks are logical or mixed. (iv) Standard SFL was weak in our setting, and the extra SFL seeds are post hoc. (v) The bound's constant is numerically vacuous; it motivates the features rather than predicting values. (vi) An earlier version of the curriculum result was misleading; Appendix~\ref{app:history} documents what changed. Appendix~\ref{sec:limitations} discusses each point.

\subsection*{AI use statement}

In this work, we used generative AI tools to generate synthetic data, develop theoretical models and the framework, formulate mathematical claims, assist in proving them, propose hypotheses, design experiments, implement methods, assist with translation, support qualitative analysis, and interpret results. We have not used generative AI tools to clean or reformat data. Additionally, we used generative AI tools to create figures, write code, draft paper sections, edit for readability, identify gaps in the literature, and suggest the paper's structure and title.
We have reviewed all AI-assisted work. The author checked every proof line by line; ran and tested all AI-written code; regenerated every reported number from the stored result files, each of which records a hash of the code that produced it; wrote each pre-registration before the corresponding runs and logged every later change to the experiments (Appendix~\ref{app:history}); and checked every citation against its published version.
We take responsibility for the final content of this work, including text, claims, or artifacts produced with the aid of generative AI.

% Acknowledgements: camera-ready only (anonymity). With [accepted], use
% \acknowledgments{...}

\bibliographystyle{abbrvnat}
\bibliography{references}

\input{checklist}

\clearpage
\appendix
\thispagestyle{empty}
\onecolumn
\aistatstitle{Predicting Level Difficulty Before Running It:\\ Reliable Labels, Layout Graders, and Cheaper Curricula\\ Supplementary Materials}

\section{Extended Related Work}
\label{app:related}



The framework of this paper sits at the meeting point of five literatures that rarely cite one another: automatic curriculum learning, successor-feature transfer, metrics on Markov decision processes, spectral graph methods in reinforcement learning, and empirical hardness modeling for combinatorial solvers, the last of which lies outside reinforcement learning altogether. We treat each literature in turn and identify which of our components are new and which are recombinations of existing ones. Table~\ref{tab:positioning} places our method against the nearest environment-scoring baselines along the two axes that matter for our contribution: whether the score is computable without executing a policy, and whether it comes with a value-theoretic bound.

\begin{table}[h]
\centering
\small
\caption{Positioning against the closest environment-scoring methods. Execution-free indicates that the score is computable without instantiating or stepping the environment.}
\label{tab:positioning}
\small
\begin{tabular}{>{\raggedright\arraybackslash}p{5.3cm}>{\raggedright\arraybackslash}p{3.9cm}>{\raggedright\arraybackslash}p{2.9cm}>{\raggedright\arraybackslash}p{3.1cm}}
\toprule
Method & Signal & Execution-free & Value-theoretic bound \\
\midrule
PLR \citep{jiang2021prioritized}          & TD error / value loss        & no  & no \\
PAIRED / ACCEL \citep{dennis2020emergent, parker2022evolving} & approximate regret & no  & minimax regret \\
TRACED \citep{cho2025traced}              & transition-prediction error  & no  & no \\
SFL \citep{rutherford2024no}              & rollout learnability $p(1{-}p)$ & no & no \\
ProCuRL \citep{tzannetos2023proximal}     & success rate                 & no  & expected improvement \\
PACE \citep{yuan2026pace}                 & parameter-update norm        & no (provisional update) & first-order \\
PIC/POIC \citep{furuta2021policy}         & policy--return MI            & no (sampled episodes) & information-theoretic \\
Task embeddings \citep{mahajan2024learning} & agent-population performance & no & information-theoretic \\
SF\&GPI \citep{barreto2018transfer}       & reward-vector distance       & \textbf{yes} & yes (shared dynamics only) \\
\textbf{This work}                        & feasibility-graph geometry   & \textbf{yes} (labels paid once) & loose; motivates features \\
\bottomrule
\end{tabular}
\end{table}

\subsection{Where Our Components Come From}
\label{sec:related}

We claim no individual component as unprecedented. The reward and dynamics split of Eq.~\ref{eq:decomposition} is implicit throughout the successor-feature literature \citep{barreto2018transfer, lehnert2020successor}; what is new is the descriptor to which the dynamics term is tied. The Laplacian descriptor has been used in reinforcement learning for two decades \citep{mahadevan2005proto, machado2017laplacian, machado2018eigenoption}, and \citet{menache2002qcut} and \citet{simsek2005identifying} motivate our bottleneck features (bridges and articulation points); the distinction we draw is one of \emph{use}, computing the spectrum across a family as a covariate rather than within one environment as a representation. A regressor over inexpensive instance features is an empirical hardness model \citep{leytonbrown2009empirical, xu2008satzilla, hutter2014algorithm}, whose vocabulary we adopt---including ``probing feature'' for our low-cost probe---rather than reinventing. Sokoban difficulty prediction from static features predates this work by fifteen years \citep{ashlock2010evolution, vankreveld2015automated, kartal2016data}, though that literature targets a human or a complete search solver, both far more layout-sensitive than a learning agent; our Karel result is a counterexample to transferring that assumption. Our reference solver follows \citet{carvalho2024reclaiming}. What we contribute is the execution-free bound connecting these components, a reliability-gated measurement protocol the curriculum literature has operated without, and results delimiting when the predictor is worth using.



\subsection{Automatic curriculum learning and unsupervised environment design}

\citet{narvekar2020curriculum} survey curriculum learning for reinforcement learning. Task-selection strategies divide by the signal they use. Learning-progress methods select tasks by measured competence gain \citep{matiisen2019teacher, portelas2020teacher}, and goal-generation methods synthesize goals at the frontier of ability \citep{florensa2018automatic}. Zone-of-proximal-development methods score tasks by proximity to the current competence of the agent. \citet{tzannetos2023proximal} give this idea a formal footing, deriving the learnability score $p(1-p)$ and showing that in two simple settings greedy maximization of that score is equivalent to maximizing expected improvement; SFL's learnability sampling builds on the same idea.

Unsupervised environment design \citep{dennis2020emergent} casts curriculum construction as a game between a level-generating teacher and a learning student, with regret as the canonical teacher objective. The line includes learned adversaries \citep{dennis2020emergent}, replay-based curation \citep{jiang2021prioritized, jiang2021replay}, evolutionary editing \citep{parker2022evolving}, multi-agent extensions \citep{samvelyan2023maestro}, and refinements of the regret objective \citep{beukman2024refining}. In parallel, the open-endedness literature co-evolves environments with their solutions under a minimal-criterion filter \citep{brant2017minimal, wang2019poet, wang2020enhanced}, and quality-diversity methods maintain archives spread across a behavior descriptor space \citep{mouret2015illuminating}. Our negative selection result (Appendix~\ref{sec:results-curriculum}) suggests that, in the domains we test, choosing diverse levels by layout features adds little over choosing them at random.

\citet{rutherford2024no} provide the most directly relevant point of comparison. They show that regret approximations correlate poorly with intuitive learnability, and that sampling high-learnability levels outperforms unsupervised environment design methods. Our work is consistent with theirs and complementary in one respect. SFL evaluates learnability by rollouts over a large sampled pool, and that rollout cost is the target of our framework. Section~\ref{sec:results-curriculum-v2} tests one composition of the two, in which the grader chooses among candidates and a few rollouts calibrate it.

\subsection{Reducing or removing the rollout cost of environment scoring}

Several recent methods address the same cost from different angles and constitute the sharpest competitors to our framework.

\citet{yuan2026pace} score a level by the squared $\ell_2$ norm of the induced policy-parameter update, presenting the method as evaluation without additional rollouts. The distinction from our setting is material. PACE performs a provisional policy update on the level, so the environment must be instantiated and stepped, and what the method removes is the separate scoring rollout rather than the execution. Our geometric descriptor never instantiates the environment. PACE is rollout-light where our descriptor is execution-free, and both are legitimate operating points. We do not claim ours dominates.

\citet{furuta2021policy} propose policy information capacity and policy-optimal information capacity, the mutual information between policy parameters and episodic return and between policy parameters and episodic optimality respectively, as environment-agnostic and algorithm-agnostic measures of task difficulty. They report higher correlation with normalized solvability than a range of alternatives, and use the measures to tune design parameters such as reward shaping without running full reinforcement learning experiments. This is the closest prior work in spirit to our difficulty-measurement goal. It differs in two respects. Both quantities are estimated by sampling policies and running episodes, so they are inexpensive relative to full training but are not computable from the specification, and both characterize a task in aggregate rather than ranking instances within a parameterized family.

\citet{mahajan2024learning} learn fixed-dimensional task embeddings from a population of agents under an information-theoretic similarity criterion, and apply them to two downstream problems that we also consider, namely predicting the performance of an agent on a held-out task and selecting tasks with desired characteristics from a candidate set. Their experiments include a Karel variant. Their embedding derives from observed agent performance across a population and therefore requires the rollouts our framework avoids. It also captures semantic structure that our static descriptor demonstrably cannot (Section~\ref{sec:results-karel}). Their result and our Karel negative result point in the same direction. Performance-derived task representations succeed in settings where specification-derived representations fail.

\citet{krsteski2026predicting} study difficulty prediction before execution across seventeen agentic benchmarks, identifying token-level entropy as a predictive signal and organizing prior work by when the difficulty signal is observed, whether after, during, or before execution. We adopt that taxonomy. Their setting places language-model agents on static benchmark tasks with textual descriptions, whereas ours places reinforcement learning policies on parameterized Markov decision processes with graph-structured specifications, and we provide a value-theoretic justification for our descriptor that has no analogue in their setting.

\citet{arxiv2507graph} propose a graph-based complexity metric for multi-agent coordination tasks, validate it against empirical difficulty, and use it to order a curriculum, a claim structure parallel to ours in a different domain. Their difficulty label is measured against a random agent and they report no reliability analysis, and we regard these as the principal methodological differences.

\subsection{Successor features and transfer under changing dynamics}

Successor features \citep{barreto2017successor, barreto2018transfer} and their universal \citep{borsa2019universal} and generalized-policy-update \citep{barreto2020fast} variants provide value prediction across tasks under a linear reward decomposition and shared dynamics. Extending them to varying dynamics is an active line of work. \citet{abdolshah2021new} treat source successor features as noisy observations of the target successor-feature function under a Gaussian process, proving convergence and a bounded modeling error under differing dynamics and rewards. \citet{malekzadeh2023uncertainty} construct a hybrid model-based successor-feature algorithm that removes the fixed-dynamics assumption and generalizes across both rewards and transition dynamics, with Kalman-filter uncertainty estimates. \citet{zhang2017deep} distinguish transfer across reward from transfer across maze structure in a navigation setting, and \citet{lehnert2020successor} clarify the relationship between successor features and model-based prediction. The latter relationship is the structural fact underlying our resolvent argument.

Relative to this line, Eq.~\ref{eq:decomposition} is not a new decomposition. The reward and dynamics split is implicit throughout the literature and we do not claim it as ours. The new element is the descriptor to which the dynamics term is tied. Every method above estimates the dynamics gap from interaction data in the target environment, whereas Theorem~\ref{prop:adjacency-control} ties that gap to a combinatorial object present in the environment program before any interaction occurs.

\subsection{Measuring distance between MDPs}

The simulation lemma \citep{kearns2002near} bounds the policy-evaluation gap between two Markov decision processes by the total-variation distance between their kernels, and in its standard $\ell_\infty$ form it is dimension-free, a point taken up in Remark~\ref{rem:ellinfty}. Bisimulation metrics \citep{ferns2004metrics, ferns2011bisimulation} define a pseudometric on states combining a reward differential with a Kantorovich distance between transition distributions, and bound value differences by that pseudometric. \citet{castro2020scalable} gives scalable on-policy variants, \citet{song2016measuring} extends the construction across separate but homogeneous processes, and \citet{tao2025generalized} establish the metric properties needed for inter-process policy-transfer bounds. Related work studies Lipschitz continuity of value in model parameters \citep{hinderer2005lipschitz, rachelson2010locality, asadi2018lipschitz, pirotta2015policy} and structural similarity measures for transfer \citep{castro2010using}.

These constructions form the correct comparison class for Theorem~\ref{prop:adjacency-control} and Appendix~\ref{app:valuegap}, and the comparison does not favor us uniformly. Bisimulation-style bounds are tighter and more general than ours. We sacrifice tightness and generality in exchange for a bound whose right-hand side is a function of the environment specification rather than of the transition kernels. In the deterministic grid families we study, Corollary~\ref{cor:common-edge} makes the trade lossless because the residual vanishes. Outside those families it does not.

\subsection{Spectral and graph-theoretic structure in reinforcement learning}

The normalized graph Laplacian of the state-transition graph has a long history in reinforcement learning. Proto-value functions \citep{mahadevan2005proto, mahadevan2007proto} use its low-order eigenvectors as bases for value-function approximation, motivated by the sensitivity of those eigenvectors to bottlenecks and symmetries. \citet{machado2017laplacian} construct eigenoptions from the same spectrum. \citet{machado2018eigenoption} connect proto-value functions to the successor representation, establishing the bridge between spectral structure and successor features that Appendix~\ref{app:valuegap} also uses, and we therefore do not claim that bridge as novel. \citet{menache2002qcut} and \citet{simsek2005identifying} use minimum cut and local graph partitioning on the transition graph to identify bottleneck subgoals, and their work motivates our bottleneck features (bridges and articulation points). \citet{jinnai2019discovering}, \citet{wu2019laplacian}, \citet{wang2021towards}, and \citet{klissarov2023deep} develop scalable approximations of the Laplacian representation. The narrow-passage intuition has an independent formalization in sampling-based motion planning, in which expansiveness and $\epsilon$-goodness are geometric parameters of the free space that bound planner running time \citep{hsu1997path, kavraki1998analysis}. To our knowledge these parameters are the closest existing precedent for the general form of Hypothesis~\ref{hyp:lipschitz}.

The distinction we draw concerns use rather than object. The literature above computes the spectrum of a single environment in order to build a representation, option set, or subgoal decomposition for solving that environment. We compute the spectrum across a family, as a covariate for predicting the realized difficulty of a learning agent, and our central empirical claim concerns the conditions under which such prediction succeeds and fails.

\subsection{Empirical hardness models and algorithm selection}

Outside reinforcement learning, predicting solver performance on an instance from inexpensively computable instance features is a mature paradigm. Empirical hardness models \citep{leytonbrown2002learning, leytonbrown2009empirical} fit regressors from instance features to solver runtime, SATzilla \citep{xu2008satzilla} uses such models to select an algorithm for each instance, and \citet{hutter2014algorithm} survey the methodology across satisfiability, mixed-integer programming, and the traveling salesman problem. Instance space analysis \citep{smithmiles2012measuring, smithmiles2014towards, munoz2020instance} extends the algorithm-selection framework of \citet{rice1976algorithm} by projecting instances into a feature space and reasoning about coverage and algorithm footprints.

Our framework is an environment-design instance of this paradigm and we adopt its vocabulary rather than reinventing it. Two correspondences hold. Our grader is an empirical hardness model. The low-cost probe of Appendix~\ref{sec:mini-hc}, which runs a short fixed program library or a severely truncated solver run and uses the resulting rewards as features, is a probing feature in the sense of \citet{hutter2014algorithm}, and we adopt that term. Our addition to this paradigm is a value-theoretic account of why a particular feature family should predict performance. Features in empirical hardness models are selected by domain intuition and validated empirically, whereas our reachability and spectral layers are the descriptor singled out by Theorem~\ref{prop:adjacency-control} and the Cheeger inequality. A further difference concerns the prediction target. Empirical hardness models target complete solvers on static instances, whereas a curriculum predictor targets a learner whose competence is non-stationary, and for that reason our difficulty label requires a reliability gate that instance-hardness labels do not.

\subsection{Difficulty estimation in procedural content generation}

Games research has estimated level difficulty from static content for over a decade, and Sokoban is one of its standard testbeds, making this literature the most direct prior art for Section~\ref{sec:results-sokoban}. \citet{ashlock2010evolution} assess the difficulty of Sokoban boards using the mean time-to-solution and failure count of an evolutionary solver as a difficulty surrogate, a label construction structurally identical to Eq.~\ref{eq:difficulty-label} in the same domain. They further report that a naive move-sequence representation produced almost no usable signal, anticipating our sparse-reward degeneracy result. \citet{kartal2016data} and \citet{vankreveld2015automated} identify inexpensively computable static board features correlated with human-perceived difficulty and use them to drive generation. A separate strand shows that simple search statistics, among them A$^\star$ closed-list length and breadth-first-search nodes expanded, correlate with perceived difficulty in mazes and Sokoban \citep{vankreveld2015automated}, and more recent work regresses solvability and push count directly with convolutional networks. Human-facing difficulty has also been modeled cognitively \citep{jarusek2010difficulty} and, on the agent side, predicted using deep reinforcement learning playtesters \citep{roohi2021predicting}.

We differ from this literature in target rather than in method, and the difference is the point of our Karel result. Difficulty estimation in procedural content generation predicts difficulty for a human solver or for a complete search solver, both far more sensitive to layout than a learning agent. Our finding that Karel tasks can carry a highly reliable difficulty label ($\rho_{\mathrm{rel}}$ up to $0.82$) that static geometry cannot predict is a counterexample to the implicit transfer of that assumption to reinforcement learning. We do not claim to outperform static-feature baselines on human difficulty, and Section~\ref{sec:limitations} records as a limitation that we do not run the search-node-count baseline this literature would treat as standard.

\subsection{Programmatic reinforcement learning}

Our reference solver operates over programmatic policies. LEAPS \citep{trivedi2021learning} introduced the latent program space for Karel, while HPRL \citep{liu2023hierarchical} and the hierarchical programmatic option framework \citep{lin2024hierarchical} compose programs hierarchically. \citet{carvalho2024reclaiming} show that hill climbing directly in the programmatic space is competitive with or superior to latent-space search, and their result is the basis of our reference solver. \citet{qiu2022programmatic} study programmatic reinforcement learning without oracles, and \citet{verma2018programmatically} and \citet{bastani2018verifiable} give the interpretability motivation for the paradigm. \citet{liu2025llmgs} introduce language-model-guided search over Karel programs and, of direct relevance to our experimental design, introduce the \textsc{PathFollow} and \textsc{WallAvoider} tasks that appear in our Karel screen. \textsc{WallAvoider} is the one Karel task whose difficulty is spatial, and it is where our grader does best on Karel (Section~\ref{sec:results-karel}).


\section{Proofs and Supporting Remarks}
\label{app:proofs}

\subsection{The Value-Gap Setting}
\label{app:valuegap}

This subsection gives the setting in which Theorem~\ref{prop:adjacency-control} is a statement about values. A stationary policy $\pi$ has value $V^\pi(\theta) = \E_{\pi, \Tdist_\theta}[\sum_{t\ge0} \gamma^t R_\theta(s_t, a_t) \mid s_0 \sim I]$.

\paragraph{What successor features already give.} The classical SF\&GPI setting \citep{barreto2017successor, barreto2018transfer} assumes shared dynamics and a linear reward decomposition $R_\theta(s,a) = \phi(s,a)^\top w_\theta$, with $\phi$ a feature map shared across the family and $w_\theta$ an environment-specific weight vector. The successor features $\psi^\pi(s,a) = \E_\pi[\sum_{t\ge0} \gamma^t \phi(s_t,a_t) \mid s_0{=}s, a_0{=}a]$ factor the action-value as $Q^\pi_\theta = \psi^{\pi\top} w_\theta$, and the GPI policy obeys $Q^*_\theta - Q^{\pi_{\text{GPI}}}_\theta \le \tfrac{2}{1-\gamma}\min_i \|w_\theta - w_i\|$, which is computable from $\theta_{\text{new}}$ alone but valid only under shared dynamics.

\paragraph{The dynamics gap.} When dynamics vary, the successor features depend on $\theta$, $\psi^\pi_\theta(s,a) = \E_{\pi,\Tdist_\theta}[\sum_{t\ge0}\gamma^t\phi(s_t,a_t) \mid s_0{=}s, a_0{=}a]$, and the value difference splits into two terms:
\begin{equation}
    V^\pi(\theta) - V^\pi(\theta')
    = \underbrace{\E_{s \sim I}\!\big[\psi^\pi_{\theta'}(s)^\top (w_\theta - w_{\theta'})\big]}_{\text{(A) reward-feature gap}}
    + \underbrace{\E_{s \sim I}\!\big[(\psi^\pi_\theta(s) - \psi^\pi_{\theta'}(s))^\top w_\theta\big]}_{\text{(B) dynamics gap}}.
\label{eq:decomposition}
\end{equation}
Term~(A) is handled by SF\&GPI. Existing bounds on term~(B) need transition kernels \citep{kearns2002near, ferns2004metrics, tao2025generalized} or target samples \citep{abdolshah2021new, malekzadeh2023uncertainty}.

\paragraph{From the transition operator to the dynamics gap.} Successor features satisfy $\psi^\pi_\theta = \phi + \gamma T^\pi_\theta \psi^\pi_\theta$. The resolvent identity and a change from the $\ell_\infty$-induced norm to the spectral norm (derivation below) give
\begin{equation}
    \big\| \psi^\pi_\theta - \psi^\pi_{\theta'} \big\|
    \;\leq\;
    \frac{\gamma \, |\mathcal V|}{(1 - \gamma)^2} \, \big\| T^\pi_\theta - T^\pi_{\theta'} \big\|_{\mathrm{op}} \, \|\phi\|,
\label{eq:resolvent-bound}
\end{equation}
and combining with Theorem~\ref{prop:adjacency-control},
\begin{equation}
    \big\| \psi^\pi_\theta - \psi^\pi_{\theta'} \big\|
    \;\leq\;
    \frac{\gamma \, |\mathcal V| \, \|\phi\|}{(1 - \gamma)^2}
    \Big(
        C_\pi \big\| A_\theta - A_{\theta'} \big\|_F
        +
        \Delta_\pi(\theta,\theta')
    \Big).
\label{eq:final-bound}
\end{equation}
The first term grows with the number of edges whose feasibility changed; the second vanishes under Corollary~\ref{cor:common-edge}. For the spectral layer, Weyl's inequality gives $|\lambda_i(L_\theta) - \lambda_i(L_{\theta'})| \le \|L_\theta - L_{\theta'}\|_{\mathrm{op}}$; for the normalized Laplacian this difference depends on the changed edges and also on the degree changes they cause, so the spectral features are stable under small edits but we do not state a constant.

\begin{hypothesis}[Geometric Lipschitz]
\label{hyp:lipschitz}
For each policy $\pi$ there exist $L_\pi, \epsilon_\pi \geq 0$ with $\|\E_{s \sim I}[\psi^\pi_\theta(s) - \psi^\pi_{\theta'}(s)]\|_2 \leq L_\pi \, d_g(\theta, \theta') + \epsilon_\pi$ for all $\theta, \theta'$, where $d_g$ is a learnable metric on the descriptor space (\ref{eq:full-descriptor}).
\end{hypothesis}

\noindent Inequality (\ref{eq:final-bound}) proves the hypothesis when $d_g$ is the Frobenius distance between adjacency matrices; for the full descriptor it remains an empirical claim, which the experiments test.

\subsection{Derivation of the Resolvent Bound (\ref{eq:resolvent-bound})}

For policy $\pi$ under dynamics $\Tdist_\theta$, the successor features satisfy the fixed-point equation
\[
    \psi^\pi_\theta \;=\; \phi + \gamma T^\pi_\theta \psi^\pi_\theta,
    \qquad\text{equivalently}\qquad
    \psi^\pi_\theta \;=\; (I - \gamma T^\pi_\theta)^{-1} \phi ,
\]
the successor-feature analogue of the Bellman equation for $Q$, with the scalar reward replaced by the vector feature $\phi$. The operator $T^\pi_\theta$ acts on feature-valued functions $f : \mathcal V \to \R^d$ by $(T^\pi_\theta f)(u) = \sum_{v} T^\pi_\theta(u,v)\, f(v)$. The resolvent identity gives, for any two environments,
\[
    \psi^\pi_\theta - \psi^\pi_{\theta'}
    \;=\;
    \gamma \, (I - \gamma T^\pi_\theta)^{-1} \, (T^\pi_\theta - T^\pi_{\theta'}) \, (I - \gamma T^\pi_{\theta'})^{-1} \phi.
\]
Each transition operator is row-stochastic, so in the $\ell_\infty$-induced (max-row-sum) norm $\|T^\pi_\theta\|_\infty = 1$ and hence $\|(I - \gamma T^\pi_\theta)^{-1}\|_\infty \le 1/(1-\gamma)$. Converting to the spectral norm on the $|\mathcal V|$-dimensional state space costs a factor $\sqrt{|\mathcal V|}$ per resolvent by norm equivalence, which yields (\ref{eq:resolvent-bound}). The factor $|\mathcal V|$ is the price of measuring in the spectral norm, in which a stochastic operator need not have unit norm.

\begin{remark}[On the looseness of the constant, and a tighter alternative]
\label{rem:ellinfty}
The constant in~(\ref{eq:resolvent-bound}) is large. At $|\mathcal V| \approx 100$, corresponding to a $10\times10$ Sokoban board, and $\gamma = 0.99$, the prefactor $\gamma|\mathcal V|/(1-\gamma)^2$ exceeds $10^{6}$, so the bound is numerically vacuous at every scale we run. We use it structurally, to certify that the dynamics gap~(B) of~(\ref{eq:decomposition}) is controlled by a feasibility-graph distance and thereby to license the descriptor, and never to track constants.

A tighter route exists. For two processes differing only in dynamics, the standard $\ell_\infty$ simulation lemma \citep{kearns2002near} gives $\|Q^\pi_\theta - Q^\pi_{\theta'}\|_\infty \le \tfrac{\gamma}{1-\gamma}\|(\Tdist_\theta - \Tdist_{\theta'})V^\pi\|_\infty$, a dimension-free bound carrying one rather than two factors of $(1-\gamma)^{-1}$. Because $\|A_\theta - A_{\theta'}\|_F^2$ upper-bounds the number of rows on which the two adjacency structures differ, an $\ell_\infty$ analogue of Theorem~\ref{prop:adjacency-control} follows at the cost of replacing the Frobenius distance with a maximum row-disagreement count. We work in the spectral norm for two reasons. It composes directly with the Weyl perturbation argument connecting $\|A_\theta - A_{\theta'}\|$ to the Laplacian spectrum used by $g_{\mathrm{spec}}$, and the Frobenius distance is the quantity our descriptor varies. A reader concerned with constants rather than with the spectral layer should prefer the $\ell_\infty$ form.
\end{remark}

\subsection{Proof of Theorem~\ref{prop:adjacency-control}}
The proof partitions the entries of $T^\pi_\theta - T^\pi_{\theta'}$ by how feasibility compares across the two environments: pairs infeasible in both contribute nothing by support consistency, pairs feasible in exactly one are charged to $\|A_\theta - A_{\theta'}\|_F$, and pairs feasible in both are collected into $\Delta_\pi$. The residual cannot be dropped in general (Remark~\ref{rem:residual-sharp}), but it vanishes in the setting we study.


\begin{proof}
Let
\[
    M := T^\pi_\theta - T^\pi_{\theta'}.
\]
We prove the estimate by splitting $M$ according to how the feasibility
of each pair $(u,v)$ compares across the two environments.

Fix arbitrary vertices $u,v\in\mathcal V$. By definition,
\begin{equation}\label{eq:pf-entry}
    M(u,v)
    =
    \sum_{a\in\mathcal A}
    \pi(a\mid u)
    \Bigl(
        \mathcal T_\theta(v\mid u,a)
        -
        \mathcal T_{\theta'}(v\mid u,a)
    \Bigr).
\end{equation}
We partition the pairs $(u,v)$ into three groups and treat each in turn.

\emph{Group 1: infeasible in both,} $A_\theta(u,v)=A_{\theta'}(u,v)=0$.
By support consistency,
\[
    \mathcal T_\theta(v\mid u,a)=0
    \quad\text{and}\quad
    \mathcal T_{\theta'}(v\mid u,a)=0
    \qquad
    \text{for every } a\in\mathcal A.
\]
Therefore every summand in~(\ref{eq:pf-entry}) is zero, and so
\begin{equation}\label{eq:pf-group1}
    M(u,v)=0 .
\end{equation}

\emph{Group 2: feasible in exactly one,} $A_\theta(u,v)\neq A_{\theta'}(u,v)$.
These are precisely the entries on which the two adjacency matrices
differ; their magnitude is bounded in the remainder of the proof.

\emph{Group 3: feasible in both,} $A_\theta(u,v)=A_{\theta'}(u,v)=1$.
Here both kernels may place mass on $(u,v)$, and nothing forces them to
place the \emph{same} mass. The corresponding entries of $M$ are
the terms collected by the shared-edge residual~(\ref{eq:residual}), so
writing $M_{\mathrm{sh}}$ for the restriction of $M$ to Group~3 we have,
by definition,
\begin{equation}\label{eq:pf-group3}
    \|M_{\mathrm{sh}}\|_F = \Delta_\pi(\theta,\theta') .
\end{equation}

By~(\ref{eq:pf-group1}) the matrix $M$ vanishes on Group~1, so its
support lies in Groups~2 and~3. Writing $M_{\mathrm{df}}$ for the
restriction of $M$ to Group~2, we obtain the decomposition
\begin{equation}\label{eq:pf-split}
    M = M_{\mathrm{df}} + M_{\mathrm{sh}},
\end{equation}
in which the two summands have disjoint supports.

It remains to bound the magnitude of $M(u,v)$ on Group~2. Define
\[
    \mathcal A_{u,v}^{\theta,\theta'}
    :=
    \Bigl\{
        a\in\mathcal A:
        \mathcal T_\theta(v\mid u,a)>0
        \text{ or }
        \mathcal T_{\theta'}(v\mid u,a)>0
    \Bigr\}.
\]
If $a\notin\mathcal A_{u,v}^{\theta,\theta'}$, then both transition
probabilities are zero and the corresponding summand vanishes. Hence
\[
    M(u,v)
    =
    \sum_{a\in\mathcal A_{u,v}^{\theta,\theta'}}
    \pi(a\mid u)
    \Bigl(
        \mathcal T_\theta(v\mid u,a)
        -
        \mathcal T_{\theta'}(v\mid u,a)
    \Bigr).
\]
Using the triangle inequality,
\[
    |M(u,v)|
    \le
    \sum_{a\in\mathcal A_{u,v}^{\theta,\theta'}}
    \pi(a\mid u)
    \left|
        \mathcal T_\theta(v\mid u,a)
        -
        \mathcal T_{\theta'}(v\mid u,a)
    \right|.
\]
Since both $\mathcal T_\theta(v\mid u,a)$ and
$\mathcal T_{\theta'}(v\mid u,a)$ are probabilities, their difference
has absolute value at most $1$. Therefore
\[
    |M(u,v)|
    \le
    \sum_{a\in\mathcal A_{u,v}^{\theta,\theta'}}
    \pi(a\mid u).
\]
By definition of $\|\pi\|_\infty$,
\[
    \pi(a\mid u)\le \|\pi\|_\infty
    \qquad
    \text{for every } u\in\mathcal V,\ a\in\mathcal A.
\]
Moreover, by the assumed action-multiplicity bound,
\[
    \bigl|\mathcal A_{u,v}^{\theta,\theta'}\bigr|
    \le d_{\max}.
\]
Thus
\begin{equation}\label{eq:pf-mag}
    |M(u,v)|
    \le
    d_{\max}\|\pi\|_\infty
    =
    C_\pi.
\end{equation}

Since $M_{\mathrm{df}}$ is supported on Group~2, where
$\bigl|A_\theta(u,v)-A_{\theta'}(u,v)\bigr|=1$, the magnitude
bound~(\ref{eq:pf-mag}) gives the entrywise estimate
\begin{equation}\label{eq:pf-entrywise}
    |M_{\mathrm{df}}(u,v)|
    \le
    C_\pi
    \bigl|A_\theta(u,v)-A_{\theta'}(u,v)\bigr|
    \qquad
    \text{for all } u,v\in\mathcal V.
\end{equation}
Squaring and summing~(\ref{eq:pf-entrywise}) over all entries gives
\[
    \|M_{\mathrm{df}}\|_F^2
    \le
    C_\pi^2
    \sum_{u,v\in\mathcal V}
    \bigl|A_\theta(u,v)-A_{\theta'}(u,v)\bigr|^2
    =
    C_\pi^2
    \|A_\theta-A_{\theta'}\|_F^2,
\]
and hence
\begin{equation}\label{eq:pf-frob}
    \|M_{\mathrm{df}}\|_F
    \le
    C_\pi
    \|A_\theta-A_{\theta'}\|_F.
\end{equation}

Finally, the spectral operator norm is bounded by the Frobenius norm, so
applying the triangle inequality to the
decomposition~(\ref{eq:pf-split}),
\[
    \|M\|_{\mathrm{op}}
    \le
    \|M_{\mathrm{df}}\|_{\mathrm{op}}
    +
    \|M_{\mathrm{sh}}\|_{\mathrm{op}}
    \le
    \|M_{\mathrm{df}}\|_F
    +
    \|M_{\mathrm{sh}}\|_F .
\]
Substituting~(\ref{eq:pf-frob}) and~(\ref{eq:pf-group3}) and recalling
$M=T^\pi_\theta-T^\pi_{\theta'}$ yields
\[
    \bigl\|T^\pi_\theta-T^\pi_{\theta'}\bigr\|_{\mathrm{op}}
    \le
    C_\pi
    \bigl\|A_\theta-A_{\theta'}\bigr\|_F
    +
    \Delta_\pi(\theta,\theta'),
\]
as claimed.
\end{proof}

\subsection{Further Corollaries and Remarks}
\noindent Inequality (\ref{eq:final-bound}) is a \emph{partial} proof of Hypothesis~\ref{hyp:lipschitz}: it establishes the bound when $d_g$ is the Frobenius distance between adjacency matrices, with $L_\pi = \gamma|\mathcal V| C_\pi \|\phi\|/(1-\gamma)^2$ and $\epsilon_\pi = \gamma|\mathcal V|\|\phi\|(1-\gamma)^{-2}\sup_{\theta,\theta'}\Delta_\pi(\theta,\theta')$. The slack is therefore not a free parameter but the shared-edge residual, the part of the dynamics gap invisible to the feasibility graph, and it vanishes in all domains we study. The hypothesis in full---covering the local and spectral layers and arbitrary stochastic policies---remains an empirical claim, and testing it is the business of Section~\ref{sec:results}.



\begin{corollary}[Graceful degradation]
\label{cor:degradation}
Suppose the two environments differ by at most $\delta$ on their shared
edges, in the sense that
\[
    \sum_{a\in\mathcal A}
    \pi(a\mid u)
    \bigl|
        \mathcal T_\theta(v\mid u,a)-\mathcal T_{\theta'}(v\mid u,a)
    \bigr|
    \;\le\;
    \delta
    \qquad
    \text{whenever } A_\theta(u,v)=A_{\theta'}(u,v)=1 .
\]
Since at most $d_{\max}$ shared edges leave each of the $|\mathcal V|$
vertices, at most $d_{\max}|\mathcal V|$ terms appear
in~(\ref{eq:residual}), and therefore
\[
    \Delta_\pi(\theta,\theta')
    \;\le\;
    \delta\sqrt{d_{\max}|\mathcal V|} .
\]
The bound~(\ref{eq:adjacency-control}) thus degrades linearly in
$\delta$: a small perturbation of the shared-edge probabilities produces
a proportionally small residual rather than invalidating the estimate,
and Corollary~\ref{cor:common-edge} is the limiting case $\delta=0$.
\end{corollary}

\begin{remark}[The residual term is necessary]
\label{rem:residual-sharp}
The residual cannot be omitted from~(\ref{eq:adjacency-control}).
Consider two environments with identical feasibility structure,
$A_\theta=A_{\theta'}$, that nonetheless assign different probabilities
to those shared edges, so that $T^\pi_\theta\neq T^\pi_{\theta'}$. Then
\[
    \|A_\theta-A_{\theta'}\|_F=0
    \qquad
    \text{but}
    \qquad
    \|T^\pi_\theta-T^\pi_{\theta'}\|_{\mathrm{op}}>0,
\]
so the adjacency term alone cannot bound the left-hand side, and the
entire discrepancy is carried by $\Delta_\pi(\theta,\theta')$. The two
terms in~(\ref{eq:adjacency-control}) therefore measure different
quantities. The first counts edges whose feasibility changed and the
second measures disagreement on edges whose feasibility did not change,
so neither term is redundant. Nothing in the statement restricts which
edges exist. The two environments may have entirely different
connectivity, and $\|A_\theta-A_{\theta'}\|_F$ measures the edges present
in one but absent in the other.
\end{remark}

\begin{remark}[The graph is unweighted]
\label{rem:weights}
The adjacency matrix $A_\theta$ records whether a move is possible rather than how probable it is, because the hardness of a navigation problem is governed by the topology of what is reachable rather than by step-by-step probabilities. A bottleneck is hard because only one route passes through it, and whether the agent traverses that passage with probability $0.7$ or $0.99$ changes neither the existence nor the narrowness of the passage. Transition probabilities govern how quickly a policy finds and exploits a path, but they do not change which paths exist.

The bound~(\ref{eq:final-bound}) reflects this structure. It charges only the moves whose status changed between $\theta$ and $\theta'$, meaning edges that opened or closed, and contributes nothing from moves feasible in both environments provided those carry identical probabilities, since such moves cancel in $T^\pi_\theta - T^\pi_{\theta'}$. When two environments share a feasible edge but assign it different step probabilities the cancellation fails, the binary adjacency alone no longer suffices, and the shared-edge residual $\Delta_\pi$ of~(\ref{eq:residual}) accounts for the difference. Whenever that residual vanishes, as it does in the deterministic grid domains we study (Corollary~\ref{cor:common-edge}), the unweighted graph captures the structural information the bound requires.
\end{remark}

\begin{remark}[Full value gap bound]
Substituting Hypothesis~\ref{hyp:lipschitz} into the decomposition~(\ref{eq:decomposition}) and bounding the reward term by the SF\&GPI argument of Section~\ref{sec:prelim},
\begin{equation}
    \big| V^\pi(\theta) - V^\pi(\theta') \big|
    \;\leq\;
    \underbrace{\|\psi^\pi_{\theta'}\|_\infty \cdot \|w_\theta - w_{\theta'}\|}_{\text{reward term, from SF\&GPI}}
    \;+\;
    \underbrace{\|w_\theta\| \cdot \big( L_\pi \, d_g(\theta, \theta') + \epsilon_\pi \big)}_{\text{dynamics term, from feasibility graph}}.
\end{equation}
The reward term uses the descriptor $w_\theta$ and is the classical SF\&GPI bound \citep{barreto2018transfer}. The dynamics term uses the geometric descriptor $g(\theta)$ of~(\ref{eq:full-descriptor}) and is the new contribution.
\end{remark}
\section{Limitations: Full Discussion}
\label{sec:limitations}

We list the limitations we consider material, in approximate order of how much each constrains the conclusions.

\textbf{The grader ties a good designer knob.} On \textsc{LavaCrossing} the grader orders a curriculum as well as grid size does, but not better ($p=0.71$ at $20$ seeds, $0.35$ at $40$, size numerically ahead). Its advantage is that it needs no knob, is not told which knob matters, and ranks levels within a size (pooled within-size Spearman $0.24$); river count, an equally simple knob, was no better than random. We tested a knob-free generator (structured mazes, Appendix~\ref{app:maze}), but found no setting where order could matter: with many shortcuts flat training already succeeded, and with fewer shortcuts neither flat training nor a simple wall-count curriculum learned. A knob-free generator with a learnable-only-with-curriculum regime remains to be tested.

\textbf{Label budget.} The grader is a learned model and needs labelled levels. On \textsc{LavaCrossing} these came from six reference agents ($3.6$M steps), more than one SFL run's scoring budget, so the grader saves compute only when reused. It beats a CNN on the raw grid up to about $1{,}000$ labels and loses slightly beyond that; with abundant labels the hand-designed features are not essential.

\textbf{Logical difficulty.} On Karel the grader is worse than a single search run. Some popular benchmarks are logical (XLand-MiniGrid) or mixed (BabyAI, Craftax), and Procgen's reactive games depend on timing a static layout cannot show.

\textbf{SFL comparison.} Standard SFL learned little in our setting ($0.07$--$0.13$), so ``matching'' it is a modest bar; the lead over cheap SFL is significant only with ten post hoc seeds. Unlocks in the curriculum controller are mostly deadline-driven, and held-out evaluation uses a stochastic policy ($4$ episodes per level).

\textbf{Statistical power.} Per-seed curriculum outcomes are bimodal (seeds end near $0$ or near $0.6$--$0.9$; pooled SD $\approx0.31$), so $20$ seeds resolve differences of about $0.25$; smaller differences are reported as not resolved, never as equal. Prediction populations are $60$ to $200$ levels outside public Sokoban, and point estimates move by about $\pm0.05$ with the forest seed and fold split.

\textbf{Looseness of the bound.} The constant in (\ref{eq:final-bound}) is numerically vacuous at realistic $|\mathcal V|$ and $\gamma$ (Remark~\ref{rem:ellinfty}); it motivates the features and nothing further. Bisimulation-style bounds \citep{ferns2004metrics, tao2025generalized} are tighter and more general; our advantage is computability from the specification alone.

\textbf{Missing baselines.} We do not run performance-derived task embeddings \citep{mahajan2024learning} or policy information capacity \citep{furuta2021policy}, both of which require rollouts, or a larger LLM judge than the $7$B model.

\textbf{Scope of the formal setting.} All results assume discrete, deterministic dynamics on a shared grid, under which Corollary~\ref{cor:common-edge} makes the residual vanish; stochastic families, continuous state spaces, and physics-based generators remain untested.

\section{Implementation Details}
\label{app:method}

\subsection{Environment Sampling and Solver Configuration}



For each task we sample a population of grids by varying the generation seed ($\theta = \text{seed}$); the task objective is held fixed and only the layout changes. Each grid is labeled by Eq.~(\ref{eq:difficulty-label}) with $R$ independent hill-climbing search seeds. The reference solver uses the configuration of \citet{trivedi2021learning}: a programmatic search space with mutation scale $\sigma$, and hill climbing with neighborhood size $k=250$ and elite count $e=2$. In Karel we use the programmatic policy space with $\sigma=0.1$ and a budget of $T=5000$ program evaluations; in MiniGrid we use the MiniGrid programmatic space with $\sigma=0.25$ and $T=2000$. We use $R=12$ search seeds for every grid throughout. Population sizes are $200$ grids for \textsc{TopOff}, the eight additional Karel tasks, and the two shaped MiniGrid foils (\textsc{RedBlueDoor}, \textsc{PutNear}); $80$ grids for \textsc{Snake}, \textsc{Harvester}, and \textsc{Seeder}; and $60$ grids for shaped \textsc{LavaGap}.

\subsection{Descriptor Feature Lists}

The geometric descriptor of Section~\ref{sec:geometric-predictor} is instantiated for grid environments as follows. The feasibility graph $G_\theta$ is built over free cells, with an edge wherever a one-step move between adjacent free cells is possible under $\theta$.

\begin{itemize}[leftmargin=*, itemsep=1pt]
\item \textbf{Local ($g_{\mathrm{loc}}$).} Grid dimensions; counts of free cells, walls, and markers; dead-end count; corridor-width and branch statistics near the start cell; goal-distance statistics.
\item \textbf{Reachability ($g_{\mathrm{reach}}$).} Start-cell eccentricity; start-to-goal distance and its excess over the Manhattan distance; degree statistics; dead-end, corridor, and branch-cell counts; articulation-point and bridge counts. (\textsc{LavaCrossing} grader: the $16$ features \texttt{num\_free}, \texttt{lava\_frac}, \texttt{mean\_degree}, \texttt{deadends}, \texttt{corridor}, \texttt{branch}, \texttt{bridges}, \texttt{articulation}, \texttt{start\_ecc}, \texttt{sg\_dist}, \texttt{sg\_excess}, $\lambda_2,\ldots,\lambda_6$.)
\item \textbf{Spectral ($g_{\mathrm{spec}}$).} The leading eigenvalues $\lambda_2, \ldots, \lambda_{k_3+1}$ of the normalized Laplacian $L_\theta$.
\item \textbf{Reward structure.} For reward-bearing tasks, a small set of task-level features such as the number and placement of target markers.
\end{itemize}

The regressor is a random forest, chosen because it needs no feature scaling and is robust at the sample sizes available for each task. All reported multivariate scores are $5$-fold cross-validated, so the regressor is scored only on grids excluded from its training folds.

\subsection{Probe Configuration}
\label{sec:mini-hc}
The grader never executes the environment and so cannot see how layout and reward interact. As a supporting experiment we use a probing feature in the sense of the algorithm-selection literature \citep{xu2008satzilla, hutter2014algorithm}: run a solver for a small fixed budget and record its behavior. A truncated-search variant runs the reference solver under a budget $T_{\mathrm{probe}} \ll T$, on the order of ten evaluations rather than thousands, scoring
\begin{equation}
    D_{\mathrm{probe}}(\theta) \;=\; 1 - r^\star_{T_{\mathrm{probe}}}(\theta).
    \label{eq:mini-hc}
\end{equation}
The probe executes the environment, but for well under a second, to \emph{classify} it rather than solve it; it is not part of the grader.

The fixed-library probe executes a small set of short hand-specified programs, among them commands to move forward, to move and place a marker, and to follow a wall while clear, together with a small number of randomly sampled programs from the domain-specific language. Each is run once and its reward, crash flag, and step count recorded. Features supplied to the difficulty regressor are the rewards attained by each library program together with summary statistics over the random programs. We also report the single most predictive probe feature on its own, which throughout is \texttt{best\_probe\_reward}, the best reward reached by any program in the library.

The truncated-search variant of Eq.~(\ref{eq:mini-hc}) starts from a random base program $r_0$, with no language model involved, and evaluates roughly ten perturbations under a budget $T_{\mathrm{probe}} \ll T$. Its score mirrors the reference label: on an easy environment a small number of perturbations reaches high reward, indicating an informative local search landscape, whereas on a hard environment the reward remains stagnant under the same effort.

\subsection{Compute and Licenses}
\label{app:compute}
All experiments ran on an Apple MacBook Pro with an M4 Pro chip. Search labels, PPO, and SFL ran on CPU (at most $8$ concurrent single-threaded jobs; $12$ for the maze experiment); the CNN baseline and the LLM judge (Qwen2.5-7B-Instruct, $4$-bit, via MLX) ran on the integrated GPU. The existing assets we use are released under open-source licenses: MiniGrid \citep{chevalierboisvert2018minigrid} under Apache~2.0, the \textsc{Boxoban} levels \citep{guez2019investigation} under Apache~2.0, the Karel environment and DSL of \citet{trivedi2021learning} under the MIT license, and Stable-Baselines3 \citep{raffin2021stable} under the MIT license.

\section{Supporting Results}
\label{app:extra}

\subsection{Karel: Seed Averaging}
\label{app:seedavg}
Seed averaging is necessary. A single hill-climbing run per grid yields a difficulty that is close to irreproducible, with rank correlation $\rho \approx 0.03$ between two independent single-run measurements of the same population: difficulty from one run measures the search's stochasticity, not a property of the environment. Averaging over $R=12$ seeds through (\ref{eq:difficulty-label}) restores a reliable label on a subset of tasks (Table~\ref{tab:karel}); \textsc{Harvester} and \textsc{Seeder} retain near-zero variance across grids and are dropped as uninformative.


\subsection{Sokoban: Cross-Check with the Program-Space Solver}
\paragraph{The label is reliable and externally valid.} We label Sokoban with the same DSL and hill-climbing solver used elsewhere and---because Sokoban states are also natural game states---additionally instantiate a budgeted stochastic best-first search over board states as an independent cross-check, under definition (\ref{eq:difficulty-label}) with a multi-box generalization of the shaping, both averaged over $R=12$ seeds on a stratified pool of $200$ levels. The state-space label clears the bar decisively at $\rho_{\mathrm{rel}} = 0.985$ and is graded rather than degenerate ($\Var(D)=0.020$, range $0$--$0.64$). Crucially it is monotone in \textsc{Boxoban}'s independent tiers (mean $D$: unfiltered $0.05$, medium $0.21$, hard $0.25$; tier-vs-$D$ Spearman $0.58$), confirming that it measures a property of the level rather than an artifact of the solver, and the program-space solver reproduces this independently, the two agreeing at the level of individual instances ($\rho_s = 0.68$, $n=200$, $p<10^{-28}$).

One residual concern is that Sokoban's tier alignment is an artifact of using a state-space search, whereas Karel and MiniGrid use program-space hill climbing. We therefore re-derive the Sokoban label with the Karel and MiniGrid machinery, using a small push and move domain-specific language searched by the same hill-climbing engine (Section~\ref{sec:difficulty-label}), under the identical difficulty definition~(\ref{eq:difficulty-label}) and the same on-target and progress shaping with weights $0.8$ and $0.2$, averaged over $R=12$ seeds on the same $200$ levels. The language augments the primitive moves with a single \texttt{solveStep} macro planning one box onto its nearest free target through a bounded single-box breadth-first search, so that a finite program can in principle clear a level; this macro is the one structural difference from the primitive Karel and MiniGrid languages and we flag it explicitly. The resulting label is again monotone in the curated tiers (mean $D$: \textsc{unfiltered} $0.09$, \textsc{medium} $0.31$, \textsc{hard} $0.31$; tier-vs-$D$ Spearman $0.54$, $p<10^{-15}$) and graded ($\Var(D)=0.031$; solved fraction $0.64/0.10/0.05$ across tiers). The two independently constructed solvers agree at the level of individual instances, with Spearman $\rho_s = 0.68$ at $n=200$ and $p<10^{-28}$ (Table~\ref{tab:sokoban-solver}, Figure~\ref{fig:sokoban-dsl-hc}), indicating that the difficulty signal is a property of the level rather than of the search formulation.

Both solvers share one limitation, namely resolution. Each separates easy from hard levels but does not distinguish \textsc{medium} from \textsc{hard}, at $0.31$ against $0.31$ in program space and $0.21$ against $0.25$ in state space. This collapse is a property of the measuring instrument rather than of the levels. Difficulty is defined as one minus the best reward reached within a \emph{fixed} search budget: an easy level is solved well within the budget and earns a low score, but once a level is hard enough that the search exhausts its budget without solving it, every harder level receives the same score regardless of its true difficulty. Both \textsc{medium} and \textsc{hard} lie past this ceiling, so the difference between them is real but falls above the resolution of the budget, much as a scale that maxes out reads the same for two objects both heavier than its limit. A larger budget would in principle separate the two tiers at proportionally greater cost, and we record this as a shared limitation rather than resolving it here.

\begin{table}[h]
\centering
\caption{Two independent reference solvers, one label. The program-space DSL$+$hill-climbing solver used for Karel/MiniGrid reproduces the tier-monotonicity of the state-space search on the same $200$ \textsc{Boxoban} levels, and the two agree per level at Spearman $\rho_s = 0.68$ ($n=200$, $p<10^{-28}$).}
\label{tab:sokoban-solver}
\begin{tabular}{lcccc}
\toprule
 & \multicolumn{3}{c}{mean $D$ by tier} & tier-vs-$D$ \\
\cmidrule(lr){2-4}
Reference solver & unfilt. & medium & hard & $\rho_s$ \\
\midrule
State-space search   & $0.05$ & $0.21$ & $0.25$ & $0.58$ \\
DSL $+$ hill climbing & $0.09$ & $0.31$ & $0.31$ & $0.54$ \\
\bottomrule
\end{tabular}
\end{table}

\begin{figure}[h]
\centering
\includegraphics[width=\linewidth]{sokoban_dsl_hc.pdf}
\caption{The Sokoban difficulty label is invariant to the search formulation. \textbf{(a)} Mean difficulty $\bar D$ per \textsc{Boxoban} tier for the two reference solvers: both are monotone from \textsc{unfiltered} to the curated tiers, and the program-space DSL and hill-climbing solver (orange) recovers the same ordering as the state-space search (blue). \textbf{(b)} Agreement between the two solvers at the level of individual instances, across all $200$ levels (Spearman $\rho_s = 0.68$); colors denote the curated tier and the dashed line is $y=x$.}
\label{fig:sokoban-dsl-hc}
\end{figure}

\subsection{MiniGrid: The Non-Geometric Foils}
\label{app:foils}
We carry only \textsc{LavaGap} through the quantitative study, but we \emph{measure} rather than assume the status of the other two tasks: \textsc{RedBlueDoor} and \textsc{PutNear}, instrumented at full protocol under a task-appropriate subgoal and distance shaping, both land on the non-geometric side, with weakly reliable labels ($\rho_{\mathrm{rel}} = 0.43$ and $0.55$, at or below \textsc{TopOff}'s bar) and geometry reaching only CV $\rho = 0.34$ and $0.39$. A pilot had suggested \textsc{RedBlueDoor} might reach $\rho_{\mathrm{rel}} \approx 0.9$; that did not survive the full population, which is precisely why the gate must be computed over the complete grid set.

\subsection{Low-Cost Probe: Full Results}
\paragraph{Where geometry fails, does anything inexpensive succeed?}
\label{sec:results-mini-hc}
Partly. Against the reliable label on full populations, one probe feature---the best reward reached by any program in the fixed library---retains a modest monotone relationship with difficulty ($\rho_s = -0.38$ on \textsc{TopOff}, $-0.31$ on \textsc{Snake}), in the expected direction since grids on which a short task-exercising program scores well tend to be easier. This clearly beats the best single geometric feature ($+0.05$, $-0.09$), confirming that the probe reads a semantic signal the static descriptor cannot. At the multivariate level, however, both families are weak (geometry $0.28$, $-0.03$; probe $0.06$, $0.15$), so the probe is a coarse indicator of the semantic axis rather than a replacement for the reference solver. The truncated-search variant (\ref{eq:mini-hc}) is informative only when the reward is shaped, recovering the reference ordering on shaped \textsc{LavaGap} at $\rho_s = +0.51$ from roughly ten evaluations but uninformative on sparse Karel rewards, since a handful of random perturbations under a sparse reward almost never earns partial credit.



We evaluate the probe against the \emph{reliable} averaged label (Eq.~\ref{eq:difficulty-label}) on the full grid populations for which that label is available, namely $200$ \textsc{TopOff} and $80$ \textsc{Snake} grids. An earlier join over $50$ grids produced optimistic multivariate scores that did not survive on the full population, and we report the full-population result here. For each task we report, against the same label, the most predictive single feature and the multivariate $5$-fold cross-validated regressor, for both the geometric descriptor and the probe (Table~\ref{tab:probe}). The single-feature reading asks whether one probe quantity tracks difficulty on its own; the multivariate reading asks whether a regressor combining all of them predicts difficulty out of sample. The probe reads the reward landscape directly, whereas the geometric descriptor sees only the layout.

\begin{table}[h]
\centering
\caption{Low-cost probe versus geometric descriptor against the reliable averaged label, full populations ($n=200$ \textsc{TopOff}, $n=80$ \textsc{Snake}). The single best probe feature, the best reward over the fixed library, captures semantic difficulty that the single best geometric feature misses. The multivariate cross-validated scores are weak for both families at full population.}
\label{tab:probe}
\begin{tabular}{lll cc}
\toprule
& \multicolumn{2}{c}{Best single feature ($\rho_s$)} & \multicolumn{2}{c}{Regressor (CV $\rho$)} \\
\cmidrule(lr){2-3}\cmidrule(lr){4-5}
Task & Geometry & Probe & Geometry & Probe \\
\midrule
\textsc{TopOff} & \texttt{num\_markers}: $+0.05$ & \texttt{best\_probe}: $\mathbf{-0.38}$ & $0.28$ & $0.06$ \\
\textsc{Snake}  & \texttt{eccentricity}: $-0.09$ & \texttt{best\_probe}: $\mathbf{-0.31}$ & $-0.03$ & $0.15$ \\
\bottomrule
\end{tabular}
\end{table}

\subsection{Selecting Levels Saves Compute, but Does Not Beat Random}
\label{sec:results-curriculum}
\paragraph{Use as a prefilter.} Given a generated batch we score each level, rank, and train only on a subset of size $k \ll N$. Prefiltering asks \emph{which} levels to train on; ordering (Section~\ref{sec:results-curriculum-v2}) asks in what sequence. A predictor can help one and not the other.

We test prefiltering first in shaped \textsc{LavaGap}, the domain where the grader predicts best. We generate a pool of $N = 150$ layouts, score each with the low-cost signals of Section~\ref{sec:geometric-predictor} with no full solve, and instantiate training on a set $S$ faithfully to our difficulty definition: hill climbing over the same program space to find the single program maximizing mean shaped reward across $S$, with cost charged in environment evaluations and evaluation on $M = 80$ held-out layouts with disjoint seeds. We compare the \textbf{full} pool; a \textbf{curated} subset of size $k$ (restricted to a target-difficulty band, then greedily covering standardized descriptor space by farthest-point selection); and a \textbf{random} subset of the same size---the ablation isolating whether the predictor, rather than merely training on fewer tasks, is responsible for any effect. To make the comparison fair we draw an \emph{ensemble} of $6$ curated and $6$ random subsets and train each with $6$ search seeds, giving $36$ runs in each condition, and sweep $k \in \{10,15,20,30,50\}$.

\begin{table}[t]
\centering
\small
\caption{Shaped \textsc{LavaGap} prefilter ($N=150$ pool, $M=80$ held-out, $36$ runs in each subset condition). Held-out mean shaped reward (mean $\pm$ s.d.); $\Delta$ is curated $-$ random; $p$ is a one-sided Mann--Whitney test of curated $>$ random; ``cost'' is the full-pool training cost divided by the subset cost. The small subset matches full-pool generalization at $2.7$--$15\times$ less compute at every $k$, but curation is not distinguishable from random selection here.}
\label{tab:curriculum}
\begin{tabular}{cccccc}
\toprule
$k$ ($k/N$) & Full & Curated-$k$ & Random-$k$ & $\Delta$ (cur$-$rnd) & $p$ \;/\; cost \\
\midrule
$10$ ($7\%$)  & $0.404 \pm 0.076$ & $0.378 \pm 0.106$ & $0.370 \pm 0.082$ & $+0.008$ & $0.47$ \;/\; $14.9\times$ \\
$15$ ($10\%$) & $0.404 \pm 0.076$ & $0.390 \pm 0.100$ & $0.370 \pm 0.072$ & $+0.020$ & $0.23$ \;/\; $9.9\times$ \\
$20$ ($13\%$) & $0.404 \pm 0.076$ & $0.382 \pm 0.083$ & $0.409 \pm 0.126$ & $-0.027$ & $0.70$ \;/\; $7.6\times$ \\
$30$ ($20\%$) & $0.404 \pm 0.076$ & $0.406 \pm 0.101$ & $0.390 \pm 0.137$ & $+0.015$ & $0.08$ \;/\; $5.1\times$ \\
$50$ ($33\%$) & $0.404 \pm 0.076$ & $0.397 \pm 0.107$ & $0.407 \pm 0.118$ & $-0.010$ & $0.61$ \;/\; $2.7\times$ \\
\bottomrule
\end{tabular}
\end{table}

At every subset size, training on $k \ll N$ layouts reaches the same held-out reward as the full pool, which attains $0.404 \pm 0.076$; every subset down to $k=10$ ($7\%$ of the pool) falls within that band while consuming $2.7$ to $14.9$ times fewer environment evaluations. A large procedurally generated pool is highly redundant, and a small representative subset suffices. \textbf{The ablation, however, is a negative result and we report it as such.} Across the whole sweep the curated subset is statistically indistinguishable from a random subset of the same size: the difference is small and changes sign ($-0.027$ to $+0.020$), and no $k$ reaches significance (smallest $p = 0.08$). An earlier single-subset comparison appeared to show a large advantage at $k=15$; it did not survive the ensemble ablation, being an artifact of comparing one fortunate curated draw against a distribution of random draws. The most plausible explanation is that the \textsc{LavaGap} descriptor space is low-dimensional and the pool dense within it, so uniform sampling already covers it about as well as explicit $k$-center selection. We expect curation to matter only when the difficulty label is genuinely graded and the descriptor space is rich enough that uniform sampling misses distinct task types; a near-degenerate label, with tasks clustered at trivially easy or unsolvable and little in between, leaves the difficulty band with no learnable middle to select and the coverage step little structure to spread over. Establishing the conditions under which curation helps is left to future work. The compute saving stands on its own: it comes from pool redundancy, which the prefilter exploits with or without descriptor-aware curation.

\paragraph{Replications.} The null replicates. On Sokoban ($200$-level pool, hill-climbing learner), band-curated, random, and full-pool training reach $0.613$, $0.602$, and $0.612$ in distribution and $0.477$, $0.492$, and $0.498$ on the hard tier. With a PPO learner on a skewed $100$-level \textsc{LavaCrossing} pool ($60\%$ near-trivial layouts; subsets of $20$; $600$k steps; $6$ draws $\times$ $2$ seeds; held-out two-river layouts), random, coverage-curated, and band-plus-coverage subsets reach $0.113$, $0.114$, and $0.082$. Ordering each subset easy-to-hard by the grader raised all of them ($0.193$, $0.161$, $0.155$), so the gain comes from ordering, not from selection.

\section{Curriculum Run History, Configuration, and Secondary Results}
\label{app:history}

\paragraph{What changed from the earlier curriculum result.} An earlier version of this work reported that a closed-form ``geometry'' order matched a reference-PPO oracle on \textsc{LavaCrossing} 9/13/17 ($0.50$ vs $0.48$, flat $0.02$, $10$ seeds). An audit found three problems. (i) The geometry bins were exactly the three grid sizes, so it was a size order. (ii) $82$ of the oracle's $90$ labels were $1.0$, so its order was a sort artifact. (iii) There was no random-order control, so ``curriculum'' and ``order'' were confounded. In addition, a run on Jun 18 scored $0.085$ for the geometry condition ($0.166$ for an adaptive variant); the harness was modified before a Jun 19 rerun that scored $0.503$, and the exact change was not recorded, but it most likely replaced an oracle-fitted score with the closed-form formula. We therefore discarded that result and re-ran the experiment from a written pre-registration with a random-order control, a separate size-order condition, and the learned grader (Section~\ref{sec:results-curriculum-v2}). An earlier small-grid run with a faulty competence gate had suggested that flat training beats curricula on 7/9/11 ($0.43$ vs $0.19$); the pre-registered rerun found no difference ($0.43$ vs $0.53$).

\paragraph{Configuration.} Table~\ref{tab:curriculum-config} lists the settings. Every result file stores a hash of the harness and scoring code. Seeds are not bit-identical across CPU architectures, so all seeds of an experiment were run on one machine.

\begin{table}[htbp]
\centering
\caption{Configuration for the curriculum and SFL experiments (Section~\ref{sec:results-curriculum-v2}).}
\label{tab:curriculum-config}
\begin{tabular}{ll}
\toprule
Component & Setting \\
\midrule
Learner & PPO \citep{schulman2017proximal} (Stable-Baselines3; \citealp{raffin2021stable}), small CNN policy, CPU \\
Budget & $1.5$M env steps; held-out evaluation at checkpoints, $4$ stochastic episodes per level \\
Task pool & \textsc{LavaCrossing}, sizes $\{9,13,17\}\times$ rivers $\{1,2,3\}$; $90$ training / $36$ held-out levels \\
Observation / actions & $7\times7\times3$ egocentric; $3$ navigation actions \\
Curriculum & $3$ bins by score; gate $0.55$; deadline $0.25$ of budget; ties broken by a seeded random key \\
Grader & random forest ($300$ trees) on $16$ layout features, fit on 7/9/11 labels from $6$ reference agents \\
SFL & refresh every $250$k steps; $4$ rollouts per scored candidate; buffer of $30$ \\
Seeds & curriculum $0$--$19$ (primary), $20$--$39$ (pre-specified secondary); SFL $0$--$9$ (primary), $10$--$19$ (post hoc) \\
\bottomrule
\end{tabular}
\end{table}

\paragraph{Secondary result: grader plus LLM.} Adding the zero-shot LLM score as one extra input to the forest, trained on the same folds, gives: \textsc{LavaGap} $0.82$ (vs $0.57$); public Sokoban $0.75$, $0.71$, $0.80$ (unchanged); Sokoban $200$ $0.30$; \textsc{LavaCrossing} $0.38$--$0.46$ (unchanged); and on the gated Karel tasks \textsc{WallAvoider} $0.56$, \textsc{TopOff} $0.41$, \textsc{DoorKey} $0.32$, \textsc{Maze} $0.32$, \textsc{Snake} $0.10$. The two signals help each other where they read different things, but the combination needs an LLM call per level.

\section{Knob-Free Curriculum Test: Structured Mazes}
\label{app:maze}

\paragraph{Design (pre-registered, with amendments made before any main run).} We sought a curriculum test on a generator without a size or river-count setting. \emph{Generator:} a randomized depth-first perfect maze on a $19\times19$ grid, start top-left and goal bottom-right, after which a uniform fraction $p\in[0,0.5]$ of the removable walls is knocked out to add loops and shortcuts. The fraction $p$ is the generator's only parameter; the \emph{wall count} condition orders by the wall count it induces. \emph{Conditions:} flat, random order, wall count, path length (breadth-first start--goal distance), and the grader, a random forest on the same layout features as Section~\ref{sec:geometric-predictor}, fit on a separate pool of $90$ mazes labelled by three reference PPO agents ($1.5$M steps each, $20$ episodes per maze; $3\%$ of labels tied, agreement between agents $\rho_s=0.35$). Everything else matches Section~\ref{sec:results-curriculum-v2}: $90$ training and $36$ held-out mazes, $1.5$M steps, the same controller, seeds $0$--$19$, one-sided Mann--Whitney tests of grader $>$ random order, wall count, and path length.

\paragraph{Choosing the setting.} Pilots ($2$ seeds, flat and wall count) were used only to choose the setting, by a rule fixed in advance: the smallest size with mean flat success at most $0.30$. Mazes with independently placed random walls were solved by flat training at every size ($0.72$--$0.95$), and a generator check showed that random walls do not create detours, so we switched to the structured generator before any main run. Among structured mazes, $15\times15$ was solved flat ($0.82$) and $19\times19$ qualified ($0.00$ and $0.59$, mean $0.295$).

\paragraph{Result.} With $20$ seeds, flat training at $19\times19$ reached at least $0.25$ on $14$ seeds, so the setting was less hard than the two-seed pilot suggested. Final held-out success: flat $0.40$ $[0.28, 0.52]$, random order $0.35$ $[0.22, 0.48]$, wall count $0.47$ $[0.35, 0.59]$, path length $0.43$ $[0.30, 0.56]$, grader $0.42$ $[0.30, 0.54]$. Grader versus random order $p=0.24$, versus wall count $p=0.77$, versus path length $p=0.61$; random order versus flat $p=0.68$; wall count versus random order $p=0.053$. By the pre-registered rule the comparison is not resolved: in a setting where flat training already works, no order separates from any other.

\paragraph{Harder settings (second pre-registered addendum).} Before any further run we fixed a follow-up: the same $19\times19$ generator with fewer shortcuts, $p\in[0,0.25]$ and then $p\in[0,0.10]$ (median start--goal path $40$ and $56$ steps, against $32$ at $p\le0.5$); pilots of $4$ flat and $2$ wall-count seeds each; and a stricter rule, taking the first setting with mean flat success at most $0.15$ \emph{and} mean wall-count success at least $0.05$, so that a curriculum could plausibly help. The grader would have reused the existing calibration. Neither setting qualified. At $p\le0.25$ flat reached $0.06, 0.04, 0.00, 0.00$ and wall count $0.00, 0.01$; at $p\le0.10$ flat reached $0.00, 0.01, 0.00, 0.00$ and wall count $0.00, 0.00$. As fixed in advance, the main run was not launched. On this generator the transition from ``solved without a curriculum'' to ``not learned even with one'' was too abrupt for an ordering comparison, and we did not search further settings.

\section{Benchmark Survey: Sources}
\label{app:survey}

Each entry gives what the generator varies, with a short quote from the cited paper or repository.
\begin{itemize}[leftmargin=*, itemsep=2pt]
\item \textbf{Procgen} \citep{cobbe2020leveraging}: a level seed sets the ``level layout\ldots, the selection of game assets, the location and spawn times of entities.'' Maze, Heist, and Chaser are generated mazes; BigFish, Leaper, StarPilot, and Plunder vary spawn timing. Mostly spatial.
\item \textbf{PLR} \citep{jiang2021prioritized}: Procgen and MiniGrid MultiRoom and ObstructedMaze; the harder ObstructedMaze setting ``requires the agent to first uncover the key from under a box.'' Spatial; ObstructedMaze mixed.
\item \textbf{PAIRED} \citep{dennis2020emergent}: ``at timestep 0 it places the agent, 1 the goal, and every step afterwards an obstacle.'' Spatial (its MuJoCo hopper is out of scope).
\item \textbf{ACCEL} \citep{parker2022evolving}: ``randomly edit the block locations (by adding or removing blocks), as well as the goal location''; also MiniHack lava and BipedalWalker terrain. Spatial; BipedalWalker out of scope.
\item \textbf{SFL} \citep{rutherford2024no}: JaxNav, MiniGrid, and XLand-MiniGrid; JaxNav ``involves low-level continuous control obstacle avoidance in addition to maze solving.'' Spatial, except XLand.
\item \textbf{XLand-MiniGrid} \citep{nikulin2024xland}: rule sets and goals vary; ``the grid itself does not change.'' Logical.
\item \textbf{Craftax} \citep{matthews2024craftax}: floors differ in ``terrain generation, enemies and required skills''; the crafting tree is fixed. Mixed.
\item \textbf{MiniHack} \citep{samvelyan2021minihack}: ``distributions of grid layouts together with monsters, objects on the floor, environment features''; navigation spatial, skill tasks logical.
\item \textbf{BabyAI} \citep{chevalier2018babyai}: ``a mission combines an instruction within an initial environment state.'' Mixed.
\item \textbf{JaxUED} \citep{coward2024jaxued}: ships ``a Maze implementation'' as its core environment. Spatial.
\item \textbf{Kinetix} \citep{matthews2025kinetix}: ``procedurally generating tens of millions of 2D physics-based tasks.'' Out of scope.
\end{itemize}

\end{document}